% Crestfaktor als Randbedingung der Tonphasenwahl
% Selbstaendig kompilierbar; laesst sich ebenso per \input als Abschnitt einbinden.
\documentclass[11pt,a4paper]{article}

\usepackage[T1]{fontenc}
\usepackage[utf8]{inputenc}
\usepackage[ngerman]{babel}
\usepackage{lmodern}
\usepackage[a4paper,margin=2.5cm]{geometry}   % \textwidth = 16 cm
\usepackage{amsmath,amssymb}
\usepackage{siunitx}
\usepackage{booktabs}
\usepackage{empheq}
\usepackage{graphicx}
\usepackage[colorlinks=true,linkcolor=black,citecolor=black,urlcolor=blue]{hyperref}

\sisetup{output-decimal-marker={,}, per-mode=symbol, separate-uncertainty=true}

\emergencystretch=1.5em   % faengt die letzten Overfull-Boxen ab

\newcommand{\crest}{\ensuremath{C}}
\newcommand{\Uw}{\ensuremath{U_W}}

\title{Der Crestfaktor als Randbedingung bei der Wahl der Tonphasen}
\author{}
\date{}

\begin{document}
\maketitle

\begin{abstract}
\noindent
Die Tonphasen der beiden AODs sind freie Parameter: Jeder Wert ist am AWG
einstellbar, und die optische Rechnung ist für jeden Wert gleichermaßen
gültig. Der Crestfaktor verbietet also keinen Phasensatz. Er begrenzt
etwas anderes, nämlich die mittlere HF-Leistung, die eine
spitzenspannungsbegrenzte Verstärkerkette bei gegebenem Phasensatz noch
linear überträgt. Damit ist er kein Term der Physik des Interferenzprofils,
sondern eine Budgetgrenze der Hardware -- und genau als solche eine
legitime Nebenbedingung der Parameterwahl. Dieser Text stellt die Kette vom
HF-Signal zum Atom her, begründet die Nebenbedingung quantitativ und grenzt
den Crestfaktor gegen die optische Spitzenintensität ab, mit der er leicht
verwechselt wird.
\end{abstract}

\section{Definition und Berechnung}

Jede AOD-Achse wird mit der Summe ihrer Töne getrieben,
\begin{equation}
  s(t) \;=\; \sum_{n=0}^{N-1} a_n \cos\!\bigl(2\pi f_n t + \varphi_n\bigr),
  \qquad
  f_n \;=\; f_\mathrm{off} + \Delta f\, \frac{n}{N-1}.
  \label{eq:rf}
\end{equation}
Darin ist $f_n$ die Frequenz des $n$-ten Tons \emph{dieser Achse}; auf der
$x$-Achse also $f_0 = \SI{100.000}{\mega\hertz}$,
$f_1 = \SI{100.185}{\mega\hertz}$, $f_2 = \SI{100.370}{\mega\hertz}$. Der
Offset $f_\mathrm{off}$ ist der unterste Ton, $\Delta f$ die Spanne vom
untersten zum obersten.

\begin{table}[htb]
  \centering
  \caption{Die Symbole dieses Textes auf einen Blick.}
  \label{tab:symbols}
  \begin{tabular}{@{}p{0.09\textwidth}p{0.38\textwidth}p{0.45\textwidth}@{}}
    \toprule
    Symbol & Bedeutung & Wert am Arbeitspunkt\\
    \midrule
    $N$ & Zahl der Töne \emph{einer} Achse & $N_x=3$, $N_y=4$ (nie 12)\\
    $n$, $m$ & Tonindex der $x$- bzw.\ $y$-Achse & $0\dots N-1$\\
    $f_n$ & Frequenz des $n$-ten Tons & \SIrange{100.000}{100.370}{\mega\hertz}\\
    $f_\mathrm{off}$ & unterster Ton & \SI{100}{\mega\hertz}\\
    $\Delta f$ & Spanne, oberster minus unterster Ton & \SI{0.37}{\mega\hertz}\\
    $f_\mathrm{c}$ & \emph{frei wählbare} Bezugsfrequenz &
      hier $\bar f$, der Mittelwert der $f_n$\\
    $f_\mathrm{env}$ & ggT der Tonabstände (Mittelungsfenster) &
      \SI{0.1850}{\mega\hertz} ($x$), \SI{0.1233}{\mega\hertz} ($y$)\\
    $f_0$, $T_0$ & Grundfrequenz/-periode der \emph{optischen} Schwebung &
      \SI{61.67}{\kilo\hertz}, \SI{16.216}{\micro\second}\\
    $a_n$ & Spannungsamplitude des $n$-ten Tons & $\sqrt{r}$-Muster\\
    $\varphi_n$ & Phase des $n$-ten Tons & die gesuchten Größen\\
    $A(t)$ & komplexe Einhüllende des HF-Signals & Gl.~\eqref{eq:envdef}\\
    \bottomrule
  \end{tabular}
\end{table}

\paragraph{Zu \boldmath$f_\mathrm{c}$.}
$f_\mathrm{c}$ ist keine physikalische Größe des Aufbaus, sondern eine reine
Rechenhilfe: Sie sagt, welche schnelle Schwingung man aus dem Signal
herausziehen möchte, um die langsame Einhüllende sichtbar zu machen
(Gl.~\eqref{eq:envdef}). Die Wahl ist beliebig, denn ein anderes
$f_\mathrm{c}'$ multipliziert $A(t)$ nur mit
$e^{-\mathrm{i}2\pi(f_\mathrm{c}'-f_\mathrm{c})t}$, einem Faktor vom Betrag
eins -- $|A(t)|$ und damit der Crestfaktor bleiben unberührt. Numerisch
geprüft für $f_\mathrm{c} = 0$, $\bar f$, $f_0$, $f_{N-1}$,
\SI{100.5}{\mega\hertz} und \SI{1}{\giga\hertz}: in allen sechs Fällen
$\max|A| = 2{,}969772$, $\mathrm{rms}|A| = 1{,}714643$,
$\crest = 2{,}449426$. Im Code steht $f_\mathrm{c} = \bar f$, weil das die
Zahlen in der Rechnung klein hält.

Nicht zu verwechseln: $f_\mathrm{off}$ ist der unterste \emph{reale} Ton,
$f_\mathrm{c}$ eine gedachte Bezugslinie. Am Arbeitspunkt liegen sie mit
\SI{100.000}{\mega\hertz} und \SI{100.185}{\mega\hertz} zufällig dicht
beieinander.

Der \emph{Crestfaktor} dieses Signals ist das Verhältnis von Spitzenwert zu
Effektivwert,
\begin{equation}
  \crest \;=\; \frac{\max_t |s(t)|}{\sqrt{\bigl\langle s(t)^2\bigr\rangle_t}} .
  \label{eq:crestdef}
\end{equation}
Die folgenden Abschnitte bauen das in kleinen Schritten auf: was der
Effektivwert ist, wie man ihn bei vielen Tönen praktisch ausrechnet, und
warum die Phasen nur auf den Zähler wirken.

\subsection{Was der Effektivwert ist, und warum gerade er}

\paragraph{Schritt 1: Effektivwert.}
Der Effektivwert (englisch \emph{rms}, root mean square) einer Spannung
$s(t)$ ist
\begin{equation}
  V_\mathrm{rms} \;=\; \sqrt{\bigl\langle s(t)^2 \bigr\rangle_t}
  \;=\; \sqrt{\frac{1}{T}\int_0^{T} s(t)^2\,\mathrm{d}t} .
  \label{eq:rmsdef}
\end{equation}
Erst quadrieren, dann mitteln, dann die Wurzel -- daher der Name. Der
gewöhnliche Mittelwert $\langle s\rangle$ wäre nutzlos, er ist bei jedem
Wechselsignal null.

Der Effektivwert ist die physikalisch richtige Mittelung, weil die Leistung
in einem Widerstand quadratisch von der Spannung abhängt:
$P(t) = s(t)^2/R$. Die mittlere Leistung ist damit
\begin{equation}
  P_\mathrm{mittel} = \frac{\langle s^2\rangle}{R} = \frac{V_\mathrm{rms}^2}{R} ,
\end{equation}
also genau Gl.~\eqref{eq:powerbudget}. Anschaulich: $V_\mathrm{rms}$ ist die
Gleichspannung, die im selben Widerstand dieselbe Wärme erzeugen würde.

\paragraph{Schritt 2: der Einzelton als Eichfall.}
Für $s(t) = a\cos(2\pi f t)$ ist die Spitze offensichtlich $a$, und wegen
$\langle\cos^2\rangle = 1/2$ ist $V_\mathrm{rms} = a/\sqrt{2}$. Also
\begin{equation}
  \crest_\text{Einzelton} \;=\; \frac{a}{a/\sqrt{2}} \;=\; \sqrt{2} \;\approx\; 1{,}414 .
\end{equation}
Numerisch bestätigt: $1{,}414213$. Das ist der kleinstmögliche Wert überhaupt
-- ein einzelner Ton ist das gleichmäßigste Signal, das es gibt.

\paragraph{Schritt 3: viele Töne -- zwei Zeitskalen.}
Bei mehreren Tönen dicht beieinander (\SI{100}{\mega\hertz} Träger,
Tonabstände von \SI{0.1}{\mega\hertz}) trennen sich zwei Zeitskalen sauber:
eine schnelle Schwingung mit $\approx\SI{100}{\mega\hertz}$ und eine langsame
\emph{Einhüllende} mit einigen Mikrosekunden. Man schreibt das als
\begin{equation}
  s(t) \;=\; \mathrm{Re}\!\left[e^{\,\mathrm{i}2\pi f_\mathrm{c}t}\,A(t)\right],
  \qquad
  A(t) = \sum_n a_n\, e^{\,\mathrm{i}\left[2\pi (f_n-f_\mathrm{c})t+\varphi_n\right]} .
  \label{eq:envdef}
\end{equation}
$|A(t)|$ ist die momentane Amplitude: Über ein paar Trägerperioden ändert sie
sich praktisch nicht, das Signal schwingt dort also \emph{lokal} wie ein
Einzelton der Amplitude $|A(t)|$.

\paragraph{Schritt 4: das Signal ist ein langsam modulierter Kosinus.}
Schreibt man die Einhüllende in Betrag und Phase, $A(t)=|A(t)|\,e^{\mathrm{i}\theta(t)}$,
so wird aus Gl.~\eqref{eq:envdef}
\begin{equation}
  s(t) \;=\; \mathrm{Re}\!\left[|A|\,e^{\,\mathrm{i}(2\pi f_\mathrm{c}t+\theta)}\right]
       \;=\; |A(t)|\,\cos\!\bigl(2\pi f_\mathrm{c}t+\theta(t)\bigr) .
  \label{eq:sascos}
\end{equation}
Das Signal \emph{ist} also ein Kosinus, dessen Amplitude $|A|$ und Phase
$\theta$ sich langsam ändern. Damit ist die Spitze sofort klar: Sie wird
erreicht, wenn $|A|$ am größten ist und der Kosinus gerade durch $\pm1$ läuft,
\begin{equation}
  \max_t|s| \;=\; \max_t|A| .
\end{equation}

\paragraph{Schritt 5: woher der Faktor \boldmath$1/2$ kommt.}
Für den Effektivwert quadriert man Gl.~\eqref{eq:sascos} und benutzt die
exakte Identität $\cos^2 x = \tfrac12\left(1+\cos 2x\right)$:
\begin{equation}
  s(t)^2 \;=\; |A|^2\cos^2\!\bigl(2\pi f_\mathrm{c}t+\theta\bigr)
  \;=\; \underbrace{\tfrac12\,|A|^2}_{\text{langsam}}
      \;+\; \underbrace{\tfrac12\,|A|^2\cos\!\bigl(4\pi f_\mathrm{c}t+2\theta\bigr)}_{\text{schnell, } \approx 2f_\mathrm{c}} .
  \label{eq:s2split}
\end{equation}
Das ist noch exakt -- punktweise gegen das direkt aufsummierte Signal
geprüft, Abweichung $3\cdot10^{-12}$.

Gemittelt wird nun über die Zeit. Der erste Summand ist die gesuchte Größe.
Der zweite schwingt mit der doppelten Trägerfrequenz: Ausmultipliziert sitzen
seine Frequenzen genau bei den Summen $f_n+f_m$, am Arbeitspunkt also
zwischen \SI{200.000}{\mega\hertz} und \SI{200.740}{\mega\hertz}. Über ein
Fenster, das lang gegen $1/(2f_\mathrm{c}) \approx \SI{5}{\nano\second}$ ist,
mittelt er weg. \textbf{Das ist die einzige Näherung im ganzen Abschnitt}, und
sie ist die Trennung der beiden Zeitskalen aus Schritt~3. Der Rest ist durch
\begin{equation}
  \left|\bigl\langle \tfrac12|A|^2\cos(4\pi f_\mathrm{c}t+2\theta)\bigr\rangle\right|
  \;\lesssim\; \frac{1}{\pi\,2f_\mathrm{c}\,T}
\end{equation}
beschränkt; am Arbeitspunkt mit $T=\SI{5.405}{\micro\second}$ sind das
$2{,}9\cdot10^{-4}$. Es bleibt also
\begin{equation}
  \bigl\langle s^2 \bigr\rangle \;=\; \tfrac12\bigl\langle |A|^2\bigr\rangle .
\end{equation}

Nachgerechnet an der $x$-Achse, volles Signal inklusive
\SI{100}{\mega\hertz}-Träger auf $4\cdot10^{7}$ Stützstellen:
\begin{center}
\begin{tabular}{@{}lr@{}}
  \toprule
  $\langle s^2\rangle$ & $1{,}470243863$\\
  $\tfrac12\langle|A|^2\rangle$ & $1{,}470000000$\\
  Restterm & $2{,}44\cdot10^{-4}$ \quad(relativ $1{,}7\cdot10^{-4}$)\\
  \midrule
  $\max|s|$ & $2{,}969765$\\
  $\max|A|$ & $2{,}969772$\\
  \bottomrule
\end{tabular}
\end{center}
Der Restterm liegt unter der Schranke. Er verschwindet sogar ganz, wenn das
Fenster zusätzlich ein ganzzahliges Vielfaches der halben Trägerperiode ist:
Mit $T = 1083/(2f_\mathrm{c}) = \SI{5.4050}{\micro\second}$ statt
\SI{5.4054}{\micro\second} bleibt nur noch $4\cdot10^{-8}$. Der winzige Rest
rührt also allein daher, dass $2f_\mathrm{c}/f_\mathrm{env} = 1083{,}08$
keine ganze Zahl ist.

Eingesetzt in die Definition ergibt sich
\begin{equation}
  \crest \;=\; \frac{\max_t|s|}{\sqrt{\langle s^2\rangle}}
  \;=\; \frac{\max_t|A|}{\sqrt{\tfrac12\langle |A|^2\rangle}}
  \;=\; \sqrt{2}\;\frac{\max_t |A(t)|}{\mathrm{rms}_t |A(t)|} ,
  \label{eq:crestenv}
\end{equation}
und der \SI{100}{\mega\hertz}-Träger muss nie abgetastet werden. Probe am
Einzelton: $|A| = a$ konstant, also $\crest = \sqrt{2}\,a/a = \sqrt{2}$.
Gegen die direkte Abtastung des vollen Signals geprüft: relative Abweichung
$\le 2\cdot10^{-4}$.

\subsection{Der Effektivwert der Einhüllenden ist phasenunabhängig}\label{sec:rmsphase}

Jetzt der Schritt, der alles Weitere erklärt. Setzt man Gl.~\eqref{eq:envdef}
in $\langle|A|^2\rangle$ ein und multipliziert aus,
\begin{equation}
  \bigl\langle |A|^2 \bigr\rangle
  = \sum_n a_n^2
  + \sum_{n\neq m} a_n a_m
    \bigl\langle e^{\,\mathrm{i}\left[2\pi(f_n-f_m)t+\varphi_n-\varphi_m\right]}\bigr\rangle ,
\end{equation}
so bleiben von der Doppelsumme nur die Diagonalterme $n=m$ übrig. Das ist
der Schritt, der eine Begründung verdient, denn er gilt \emph{nicht} über ein
beliebiges Fenster.

\paragraph{Wann genau der Mischterm null ist.}
Jeder Mischterm ist ein Zeiger, der mit $\delta\! f = f_n-f_m$ umläuft. Sein
Mittel über ein Fenster $[0,T]$ lässt sich exakt ausrechnen:
\begin{equation}
  \bigl\langle e^{\,\mathrm{i}2\pi \delta\! f\, t}\bigr\rangle_{[0,T]}
  = \frac{1}{T}\int_0^{T} e^{\,\mathrm{i}2\pi \delta\! f\, t}\,\mathrm{d}t
  = \frac{e^{\,\mathrm{i}2\pi \delta\! f\, T}-1}{\mathrm{i}\,2\pi \delta\! f\, T} .
  \label{eq:phasoravg}
\end{equation}
Der Zähler verschwindet genau dann, wenn $\delta\! f\,T$ eine ganze Zahl ist
-- wenn der Zeiger also \emph{volle} Umläufe macht. Anschaulich: Ein voller
Umlauf mittelt zum Kreismittelpunkt, und wo der Zeiger dabei startet, ist
gleichgültig. \textbf{Genau deshalb fällt auch die konstante Phasendifferenz
$\varphi_n-\varphi_m$ heraus} -- sie dreht nur den Startwinkel.

Und die Bedingung ist erfüllt, weil $f_\mathrm{env}$ als größter gemeinsamer
Teiler der Tonabstände definiert ist: Jede Differenz $f_n-f_m$ ist ein
ganzzahliges Vielfaches davon. Auf der $x$-Achse etwa
\begin{equation}
  f_1-f_0 = \SI{0.185}{\mega\hertz} = 1\cdot f_\mathrm{env},\quad
  f_2-f_0 = \SI{0.370}{\mega\hertz} = 2\cdot f_\mathrm{env},\quad
  f_2-f_1 = \SI{0.185}{\mega\hertz} = 1\cdot f_\mathrm{env},
\end{equation}
also $\delta\! f\,T \in \{1,2\}$ für $T=1/f_\mathrm{env}$ und damit
Gl.~\eqref{eq:phasoravg} exakt null.

\paragraph{Mit falschem Fenster gilt nichts davon.}
Nimmt man stattdessen $T=1/\Delta f = \SI{2.703}{\micro\second}$, so wird
$\delta\! f\,T = \tfrac12$ für die beiden Nachbarpaare -- ein \emph{halber}
Umlauf, und Gl.~\eqref{eq:phasoravg} liefert $\mp0{,}6366\,\mathrm{i}$ statt
null. Dann überlebt die Doppelsumme, und mit ihr die Phasenabhängigkeit:
\begin{center}
\begin{tabular}{lcc}
  \toprule
  $\varphi_x$ & $\mathrm{rms}|A|$, $T=1/f_\mathrm{env}$ & $\mathrm{rms}|A|$, $T=1/\Delta f$\\
  \midrule
  $[0;\,0;\,0]^\circ$        & $1{,}714643$ & $1{,}714646$\\
  $[0;\,12{,}5;\,25]^\circ$  & $1{,}714643$ & $1{,}548283$\\
  $[0;\,200;\,80]^\circ$     & $1{,}714643$ & $2{,}110658$\\
  \bottomrule
\end{tabular}
\end{center}
Links steht dreimal $\sqrt{\sum a_n^2}=1{,}714643$, rechts dreimal etwas
anderes. Das ist derselbe Fehler wie im Abschnitt über das Mittelungsfenster,
nur von der anderen Seite betrachtet -- und er erklärt ihn: Mit dem falschen
Fenster hängt schon der Nenner von den Phasen ab, und der Crestfaktor misst
nicht mehr, was er messen soll.

\medskip
\noindent Mit dem richtigen Fenster bleibt also
\begin{equation}
  \mathrm{rms}|A| = \sqrt{\textstyle\sum_n a_n^2}
  \qquad\text{unabhängig von allen } \varphi_n ,
\end{equation}
numerisch bestätigt (drei völlig verschiedene Phasensätze, identischer
Effektivwert auf sechs Stellen: $1{,}714643$). Damit vereinfacht sich
Gl.~\eqref{eq:crestenv} zu
\begin{empheq}[box=\fbox]{equation}
  \crest \;=\; \sqrt{2}\;\frac{\max_t |A(t)|}{\sqrt{\sum_n a_n^2}} .
  \label{eq:crestsimple}
\end{empheq}
\textbf{Die Phasen beeinflussen den Crestfaktor ausschließlich über
$\max_t|A|$.} Sie können den Effektivwert nicht ändern, nur die Spitze --
nämlich dadurch, wie gut sich die Zeiger zu irgendeinem Zeitpunkt ausrichten
lassen. Am Arbeitspunkt:
\begin{center}
\begin{tabular}{lccc}
  \toprule
  & $\sqrt{\sum a_n^2}$ & $\max|A|$ & $\crest$\\
  \midrule
  $x$-Achse, $\varphi_x=[0;\,12{,}5;\,25]^\circ$ & $1{,}7146$ & $2{,}9698$ & $2{,}4494$\\
  $y$-Achse, $\varphi_y=[0;\,98;\,16;\,114]^\circ$ & $2{,}0785$ & $3{,}2279$ & $2{,}1963$\\
  \bottomrule
\end{tabular}
\end{center}

\paragraph{Die beiden Schranken.}
Daraus folgen sofort beide Grenzen:
\begin{align}
  \max|A| \;&\le\; \sum_n a_n
  &&\text{(Dreiecksungleichung; Gleichheit, wenn alle Zeiger ausgerichtet sind)},\\
  \max|A| \;&\ge\; \mathrm{rms}|A|
  &&\text{(ein Maximum liegt nie unter dem Effektivwert)} .
\end{align}
Eingesetzt in Gl.~\eqref{eq:crestsimple} und für gleiche Amplituden
($a_n = a$, $\sum a_n = Na$, $\sqrt{\sum a_n^2} = \sqrt{N}a$):
\begin{equation}
  \sqrt{2} \;\le\; \crest \;\le\; \sqrt{2}\,\frac{Na}{\sqrt{N}a} \;=\; \sqrt{2N} .
  \label{eq:sqrt2n}
\end{equation}
Numerisch bestätigt: $N=3\rightarrow 2{,}4495$, $N=4\rightarrow 2{,}8284$
(und, zur Einordnung, $N=13\rightarrow 5{,}0990$); am Arbeitspunkt mit
$r_x=0{,}97$ ergibt sich $\sum a_n = 2{,}970$ und
$\sqrt{\sum a_n^2}=1{,}715$, also $\crest_{\max}=2{,}4494$.

\paragraph{Amplitudenkonvention.}
Die Gewichte $r$ der Simulation sind \emph{Intensitäts}gewichte, also
HF-\emph{Leistungs}verhältnisse. In Gl.~\eqref{eq:rf} stehen Spannungen,
folglich $a_n \propto \sqrt{r}$.

\paragraph{Rezept.}
\begin{enumerate}\itemsep0.15em
  \item Tonfrequenzen $f_n$, Spannungsamplituden $a_n=\sqrt{r_n}$ und Phasen
        $\varphi_n$ zusammenstellen.
  \item $f_\mathrm{env}$ als größten gemeinsamen Teiler der Tonabstände
        bestimmen und ein Raster über \emph{eine} Periode $1/f_\mathrm{env}$ legen.
  \item $A(t)$ nach Gl.~\eqref{eq:envdef} auswerten ($f_\mathrm{c}$ beliebig,
        hier der Mittelwert -- eine gemeinsame Frequenzverschiebung ändert
        $|A|$ nicht).
  \item $\crest = \sqrt{2}\,\max|A| \big/ \sqrt{\sum a_n^2}$.
\end{enumerate}
Im Code ist das \texttt{CrestBasis} in Abschnitt~5 von
\texttt{kern/beating\_physik.py}: Die Matrix
$a_n e^{\mathrm{i}2\pi(f_n-f_\mathrm{c})t_j}$ wird einmal vorberechnet,
danach kostet eine Auswertung nur noch ein Matrix-Vektor-Produkt mit
$e^{\mathrm{i}\varphi}$.

\paragraph{Der Crestfaktor ist eine Größe \emph{je Achse}.}
Das ist die häufigste Verwechslung, deshalb explizit: Es gibt in diesem
Aufbau \emph{kein} Zwölfton-Signal. Der AWG treibt zwei getrennte Kanäle mit
\begin{align}
  N_x &= 3 \text{ Tönen bei } [\,100{,}000;\;100{,}185;\;100{,}370\,]\,\si{\mega\hertz},\\
  N_y &= 4 \text{ Tönen bei } [\,100{,}0000;\;100{,}1233;\;100{,}2467;\;100{,}3700\,]\,\si{\mega\hertz}.
\end{align}
Jeder Kanal hat seinen eigenen Verstärker, seinen eigenen AOD und damit
seinen eigenen Crestfaktor. Die zwölf Spots entstehen erst \emph{optisch}:
Jeder Spot ist nach Gl.~\eqref{eq:spotphase} das Produkt eines $x$-Tons mit
einem $y$-Ton, $3\times4 = 12$. Die HF-Ketten sehen davon nichts. In
Gl.~\eqref{eq:sqrt2n} ist $N$ deshalb $3$ bzw.\ $4$ -- niemals $12$.

\paragraph{Mittelungsfenster.}
Der Effektivwert ist über \emph{eine Periode der Einhüllenden der jeweiligen
Achse} zu bilden, $1/f_\mathrm{env}$ mit $f_\mathrm{env}$ dem größten
gemeinsamen Teiler der Tonabstände dieser Achse -- nicht über $1/\Delta f$
mit der Spanne $\Delta f$. Am Arbeitspunkt:
\begin{align}
  \text{$x$-Achse:}\quad & f_\mathrm{env} = \SI{0.1850}{\mega\hertz}
    && \rightarrow\; \text{Fenster } \SI{5.405}{\micro\second},\\
  \text{$y$-Achse:}\quad & f_\mathrm{env} = \SI{0.1233}{\mega\hertz}
    && \rightarrow\; \text{Fenster } \SI{8.108}{\micro\second},
\end{align}
gegenüber $1/\Delta f = \SI{2.703}{\micro\second}$ bei Verwendung der Spanne.
Ein zu kurzes Fenster erwischt einen unrepräsentativen Ausschnitt der
Einhüllenden; der Fehler ist erheblich und nicht einmal
vorzeichenbeständig:
\begin{center}
\begin{tabular}{lccc}
  \toprule
  & richtig ($1/f_\mathrm{env}$) & falsch ($1/\Delta f$) & Fehler\\
  \midrule
  $\crest_x$ & $2{,}4494$ & $2{,}6699$ & $+\SI{9.0}{\percent}$\\
  $\crest_y$ & $2{,}1963$ & $3{,}3679$ & $+\SI{53.3}{\percent}$\\
  \bottomrule
\end{tabular}
\end{center}
(In einer anderen Konfiguration mit 13 Tönen auf einer Achse liefert derselbe
Fehler $\crest = 2{,}19$ statt korrekt $5{,}099$ -- dort liegt das zu kurze
Fenster vollständig im Rephasierungsmaximum, der Effektivwert kommt zu groß
und der Crestfaktor zu klein heraus. Die Richtung des Fehlers hängt davon ab,
welchen Ausschnitt man erwischt.)

\paragraph{Drei Perioden, die nicht dasselbe sind.}
\begin{center}
\begin{tabular}{@{}lll@{}}
  \toprule
  HF-Einhüllende $x$ & \SI{5.405}{\micro\second} & Mittelungsfenster für $\crest_x$\\
  HF-Einhüllende $y$ & \SI{8.108}{\micro\second} & Mittelungsfenster für $\crest_y$\\
  optische Schwebung $T_0$ & \SI{16.216}{\micro\second} & Intensitätsmuster, beide Achsen zusammen\\
  \bottomrule
\end{tabular}
\end{center}
$T_0$ ist das kleinste gemeinsame Vielfache der beiden Achsenperioden
($16{,}216 = 3\cdot5{,}405 = 2\cdot8{,}108$). Der Trigger $t_0$ und alle
optischen Größen beziehen sich auf $T_0$, der Crestfaktor dagegen auf die
Periode seiner eigenen Achse.

\paragraph{Amplitudenkonvention.}
Die Gewichte $r$ der Simulation sind \emph{Intensitäts}gewichte, also
HF-\emph{Leistungs}verhältnisse. In Gl.~\eqref{eq:rf} stehen Spannungen,
folglich $a_n \propto \sqrt{r}$.

\section{Die Kopplung: eine Phase, zwei Wirkungen}

Der AOD überträgt Frequenz \emph{und} Phase des akustischen Gitters auf das
gebeugte Licht. Die $\varphi_n$ aus Gl.~\eqref{eq:rf} sind deshalb dieselben
Größen, die im optischen Feld die Interferenz bestimmen. Ein Spot wird je
einmal in $x$ und in $y$ gebeugt, sein Feld trägt also das Produkt beider
Spannungen und die Summe beider Phasen,
\begin{equation}
  E_{nm} \;\propto\; V_x(n)\,V_y(m),
  \qquad
  \varphi_{nm} \;=\; \varphi_x(n) + \varphi_y(m).
  \label{eq:spotphase}
\end{equation}
Hierin liegt der gesamte Zusammenhang: \emph{ein} Satz Zahlen legt
gleichzeitig das Interferenzmuster am Atom und die Kurvenform am Verstärker
fest. Beide lassen sich nicht unabhängig voneinander einstellen.

Wichtig für die Einordnung: Der Crestfaktor ist kein Term der Feldrechnung.
Diese setzt voraus, dass jeder Ton genau seine Amplitude erhält, dass die
Beugung linear ist und dass keine Intermodulation auftritt. Der Crestfaktor
prüft, ob diese Voraussetzung in der realen Kette Bestand hat. Er ist eine
Gültigkeitsbedingung, keine Zutat zum Modell.

\section{Warum daraus eine Nebenbedingung wird}

Der Ausgangspunkt ist die Definition Gl.~\eqref{eq:crestdef}, nach Spitzen-
und Effektivwert aufgelöst,
\begin{equation}
  V_\mathrm{peak} \;=\; \crest \cdot V_\mathrm{rms} .
  \label{eq:peakrms}
\end{equation}
Ist der Verstärker durch seine \emph{Spitzenspannung} begrenzt -- er geht in
Kompression bzw.\ Begrenzung, nicht in thermische Sättigung --, so lautet die
Hardwarebedingung $V_\mathrm{peak}\le V_{\max}$, nach
Gl.~\eqref{eq:peakrms} also $V_\mathrm{rms}\le V_{\max}/\crest$. Die an einer
reellen Last $R$ umgesetzte \emph{mittlere} Leistung ist
$P_\mathrm{mittel}=V_\mathrm{rms}^2/R$, und damit
\begin{equation}
  P_\mathrm{mittel} \;=\; \frac{V_\mathrm{rms}^{2}}{R}
  \;\le\; \frac{V_{\max}^2}{R\,\crest^{2}} .
  \label{eq:powerbudget}
\end{equation}
Mehr als die Definition des Crestfaktors und $P=V_\mathrm{rms}^2/R$ steckt
nicht darin. Gleichwertig in Leistungen ausgedrückt: Die
Spitzenhüllkurvenleistung ist $P_\mathrm{peak}=\crest^{2}P_\mathrm{mittel}$,
und verträgt die Kette $P_\mathrm{peak}\le P_\mathrm{sat}$, so folgt
$P_\mathrm{mittel}\le P_\mathrm{sat}/\crest^{2}$. Die Systemimpedanz $R$
kürzt sich aus jedem Vergleich zweier Phasensätze heraus; nur das Verhältnis
$1/\crest^{2}$ bleibt stehen.

Die Annahme dahinter ist eine \emph{harte} Spannungsgrenze mit linearem
Verhalten bis dorthin. Reale Verstärker komprimieren allmählich, und der
angegebene $\SI{1}{\deci\bel}$-Kompressionspunkt bezieht sich üblicherweise
auf einen Einzelton im Dauerstrich; für ein Mehrtonsignal ist die
Spitzenhüllkurvenleistung die maßgebliche Größe. Gl.~\eqref{eq:powerbudget}
ist damit eine obere Abschätzung, kein Ersatz für eine gemessene
Kompressionskurve.
Die gebeugte optische Leistung folgt im linearen Bereich der mittleren
HF-Leistung, und nach Gl.~\eqref{eq:spotphase} wird jeder Spot zweimal
gebeugt. Für die Spotintensität gilt daher
\begin{equation}
  I \;\propto\; P_x\,P_y \;\propto\; \frac{1}{\crest_x^{2}\,\crest_y^{2}} .
  \label{eq:intensity}
\end{equation}
Ein hoher Crestfaktor kostet also Licht, ohne dass die Simulation davon etwas
merkt. Wird die Grenze überschritten, treten zusätzlich
Intermodulationsprodukte auf: neue Frequenzen bei Summen und Differenzen der
Töne, also Geisterspots, und die Amplitudenverhältnisse verschieben sich --
das eingestellte Profil ist dann nicht mehr das gerechnete.

\section{Der Zielkonflikt}

Die für das Atom maßgebliche Größe ist die Gleichmäßigkeit der
aufgesammelten Pulsfläche über die Aufenthaltsverteilung $W(r)$ des Atoms,
\begin{equation}
  \Uw \;=\; \frac{\sigma_W(\theta)}{\langle\theta\rangle_W},
  \qquad
  \theta(r) \;=\; \int_{t_0}^{t_0+T_p} \Omega(r,t)\,\mathrm{d}t .
  \label{eq:uw}
\end{equation}
Ein seitlich versetzter Interferenzfleck erzeugt einen linearen Gradienten
quer durch die Atomwolke; ein Gradient ist genau das, was $\Uw$ misst.
„Gleichmäßig beleuchtet`` und „mittig auf dem Atom`` sind damit dieselbe
Bedingung.

Der Konflikt entsteht dadurch, dass der Zeitpunkt, zu dem der Fleck sauber,
symmetrisch und mittig steht, derselbe ist, zu dem die Töne rephasieren --
und Rephasieren treibt den Crestfaktor nach Gl.~\eqref{eq:sqrt2n} an sein
Maximum. Gleichmäßige Beleuchtung und schonende HF-Kette stehen also direkt
gegeneinander. Die Crest-Schranke ist der Punkt, an dem man sich auf dieser
Kurve platziert.

\section{Das Zentrierungskriterium}\label{sec:zentrierung}

\subsection{Warum „alles ist symmetrisch`` nicht genügt}

Der Aufbau ist in mehrfacher Hinsicht spiegelsymmetrisch. Legt man den
Ursprung in den Profilmittelpunkt und schreibt für die Spiegelung der Indizes
\begin{equation}
  \bar n \;\equiv\; N-1-n , \qquad \bar m \;\equiv\; M-1-m ,
  \label{eq:barindex}
\end{equation}
so gelten drei voneinander unabhängige Aussagen:
\begin{align}
  \text{Positionen:}&&  c_{\bar n\bar m} &= -\,c_{nm} , \label{eq:symgeo}\\
  \text{Amplituden:}&&  A_{\bar n} &= A_n , \qquad B_{\bar m} = B_m
     \qquad\text{(Muster } [r,1,\dots,1,r]\text{)} , \label{eq:symamp}\\
  \text{Frequenzen:}&&  f_{\bar n} &= F - f_n , \qquad g_{\bar m} = G - g_m .
     \label{eq:symfreq}
\end{align}

\paragraph{Nachweis von Gl.~\eqref{eq:symfreq}.}
Aus $f_n = f_\mathrm{off}+\Delta f\,\dfrac{n}{N-1}$ folgt
\begin{equation}
  f_{\bar n} \;=\; f_\mathrm{off}+\Delta f\,\frac{N-1-n}{N-1}
             \;=\; f_\mathrm{off}+\Delta f-\Delta f\,\frac{n}{N-1} ,
\end{equation}
und damit für die Summe
\begin{equation}
  f_n + f_{\bar n}
  \;=\; \Bigl(f_\mathrm{off}+\Delta f\,\tfrac{n}{N-1}\Bigr)
      + \Bigl(f_\mathrm{off}+\Delta f-\Delta f\,\tfrac{n}{N-1}\Bigr)
  \;=\; 2f_\mathrm{off}+\Delta f .
  \label{eq:sumproof}
\end{equation}
Der Index $n$ fällt heraus, die Summe ist also für \emph{jedes} Spiegelpaar
dieselbe. Ihr Wert ist zugleich der des äußersten Paares,
\begin{equation}
  F \;\equiv\; f_0 + f_{N-1}
    \;=\; f_\mathrm{off} + \bigl(f_\mathrm{off}+\Delta f\bigr)
    \;=\; 2f_\mathrm{off}+\Delta f ,
\end{equation}
womit $f_{\bar n}=F-f_n$ gezeigt ist; $G$ für die $y$-Achse analog. Am
Arbeitspunkt ($f_\mathrm{off}=\SI{100}{\mega\hertz}$,
$\Delta f=\SI{0.37}{\mega\hertz}$) ist
$F=G=\SI{200.37}{\mega\hertz}$, auf beiden Achsen und für jedes Paar
numerisch bestätigt.

Wichtig ist, was Gl.~\eqref{eq:symfreq} \emph{nicht} sagt: $F$ ist keine
Frequenz des Signals, sondern die doppelte Mittenfrequenz. Die Aussage lautet,
dass die Töne \emph{paarweise symmetrisch um die Bandmitte} $F/2$ liegen -- der
Spiegelpartner ist genauso weit oberhalb der Mitte, wie der Ton unterhalb
liegt.

Der springende Punkt ist Gl.~\eqref{eq:symfreq}: Die Frequenz spiegelt
\emph{um die Mittenfrequenz}, sie kippt also relativ dazu ihr Vorzeichen. Der
Zeitfaktor $e^{\,\mathrm{i}2\pi f_n t}$ geht dabei in
$e^{\,\mathrm{i}2\pi F t}\,e^{-\mathrm{i}2\pi f_n t}$ über, also bis auf einen
gemeinsamen Faktor in sein \emph{Konjugiertes}. Die geometrische Spiegelung
ist damit keine gewöhnliche Spiegelung, sondern eine \textbf{konjugierende}
Operation. Genau deshalb reicht es nicht, dass „alles symmetrisch`` ist: Die
Phasen müssen die Konjugation \emph{mitmachen}, sonst laufen die
gespiegelten Anteile gegeneinander.

\subsection{Herleitung}

Der Ursprung liege im Profilmittelpunkt. Das Feld ist
\begin{equation}
  E(r,t) \;=\; \sum_{n=0}^{N-1}\sum_{m=0}^{M-1} A_n B_m\, u(r-c_{nm})\,
  e^{\,\mathrm{i}\left[2\pi (f_n+g_m)t \,+\, \varphi_n+\psi_m\right]},
  \label{eq:field}
\end{equation}
mit reellem, radialsymmetrischem Einzelspotprofil $u$, also $u(-\rho)=u(\rho)$.
Die Indexspiegelung $\bar n$, $\bar m$ ist in Gl.~\eqref{eq:barindex} definiert.

\paragraph{Schritt 1: Ort spiegeln.}
Einsetzen von $-r$ in Gl.~\eqref{eq:field} ändert nur das Argument von $u$:
\begin{equation}
  E(-r,t) \;=\; \sum_{n,m} A_n B_m\, u(-r-c_{nm})\,
  e^{\,\mathrm{i}\left[2\pi (f_n+g_m)t + \varphi_n+\psi_m\right]}.
  \label{eq:step1}
\end{equation}

\paragraph{Schritt 2: Summationsindizes umbenennen.}
Dies ist ein reiner Variablenwechsel, noch ohne Physik. Da
$n\mapsto\bar n$ eine Bijektion von $\{0,\dots,N-1\}$ auf sich ist, darf in
Gl.~\eqref{eq:step1} $n$ durch $\bar p$ und $m$ durch $\bar q$ ersetzt werden;
$p$ und $q$ durchlaufen dann denselben Bereich:
\begin{equation}
  E(-r,t) \;=\; \sum_{p,q} A_{\bar p} B_{\bar q}\, u(-r-c_{\bar p\bar q})\,
  e^{\,\mathrm{i}\left[2\pi (f_{\bar p}+g_{\bar q})t
  + \varphi_{\bar p}+\psi_{\bar q}\right]}.
  \label{eq:step2}
\end{equation}
Jetzt tragen \emph{alle} Faktoren gespiegelte Indizes -- das ist der Punkt, an
dem die Symmetrien angesetzt werden können.

\paragraph{Schritt 3: die drei Symmetrien einsetzen.}
Auf Gl.~\eqref{eq:step2} angewandt, jeweils mit dem laufenden Index $p$
bzw.\ $q$:
\begin{align}
  A_{\bar p} B_{\bar q} &= A_p B_q
  && \text{nach Gl.~\eqref{eq:symamp},}\\[2pt]
  u(-r-c_{\bar p\bar q}) &= u(-r+c_{pq}) = u\bigl(-(r-c_{pq})\bigr) = u(r-c_{pq})
  && \text{nach Gl.~\eqref{eq:symgeo} und } u(-\rho)=u(\rho),\\[2pt]
  f_{\bar p}+g_{\bar q} &= (F+G) - (f_p+g_q)
  && \text{nach Gl.~\eqref{eq:symfreq}.}
\end{align}
Die Phasen bleiben dabei unberührt -- über sie ist noch nichts vorausgesetzt,
sie behalten ihre gespiegelten Indizes. Damit wird aus Gl.~\eqref{eq:step2}
\begin{equation}
  E(-r,t) \;=\; e^{\,\mathrm{i}2\pi (F+G) t}
  \sum_{p,q} A_p B_q\, u(r-c_{pq})\,
  e^{-\mathrm{i}2\pi (f_p+g_q)t}\,
  e^{\,\mathrm{i}\left(\varphi_{\bar p}+\psi_{\bar q}\right)} .
  \label{eq:step3}
\end{equation}
Das Vorzeichen im Zeitfaktor ist gekippt: Der gespiegelte Spot läuft zeitlich
rückwärts.

\paragraph{Schritt 4: mit dem konjugierten Feld vergleichen.}
Das konjugierte Feld am ursprünglichen Ort lautet, aus Gl.~\eqref{eq:field}
durch komplexe Konjugation und mit demselben Summationsbuchstaben,
\begin{equation}
  E^{*}(r,t) \;=\; \sum_{p,q} A_p B_q\, u(r-c_{pq})\,
  e^{-\mathrm{i}2\pi (f_p+g_q)t}\,
  e^{-\mathrm{i}\left(\varphi_{p}+\psi_{q}\right)} .
  \label{eq:step4}
\end{equation}
Gl.~\eqref{eq:step3} und Gl.~\eqref{eq:step4} haben dieselben Basisfunktionen
$A_p B_q\,u(r-c_{pq})\,e^{-\mathrm{i}2\pi(f_p+g_q)t}$ und unterscheiden sich
ausschließlich im konstanten Phasenfaktor jedes Terms. Sie sind daher
proportional zueinander, sobald
\begin{equation}
  e^{\,\mathrm{i}\left(\varphi_{\bar p}+\psi_{\bar q}\right)}
  \;=\; e^{\,\mathrm{i}\Theta}\,
  e^{-\mathrm{i}\left(\varphi_{p}+\psi_{q}\right)}
  \qquad\text{für alle } p,q
\end{equation}
mit einer einzigen, von $p$ und $q$ unabhängigen Konstanten $\Theta$, also
\begin{equation}
  \underbrace{\bigl(\varphi_p + \varphi_{N-1-p}\bigr)}_{\text{hängt nur von } p\text{ ab}}
  \;+\;
  \underbrace{\bigl(\psi_q + \psi_{M-1-q}\bigr)}_{\text{hängt nur von } q\text{ ab}}
  \;=\; \Theta \pmod{2\pi}.
  \label{eq:separates}
\end{equation}
Eine Summe aus einem reinen $p$- und einem reinen $q$-Term ist genau dann für
alle Indexpaare konstant, wenn jeder Summand für sich konstant ist. Damit
folgt das Kriterium, und zwar \emph{je Achse getrennt}:
\begin{empheq}[box=\fbox]{equation}
  \varphi_n + \varphi_{N-1-n} \;=\; c_x ,
  \qquad
  \psi_m + \psi_{M-1-m} \;=\; c_y ,
  \qquad c_x + c_y = \Theta .
  \label{eq:centering}
\end{empheq}
Ist es erfüllt, so ist nach Gl.~\eqref{eq:step3} und Gl.~\eqref{eq:step4}
\begin{equation}
  E(-r,t) \;=\; \underbrace{e^{\,\mathrm{i}\left[2\pi (F+G)t+\Theta\right]}}_{\textstyle e^{\,\mathrm{i}\alpha(t)}}\;E^{*}(r,t)
  \quad\Longrightarrow\quad
  I(-r,t) = I(r,t)
  \quad\Longrightarrow\quad
  \nabla\theta\big|_{r=0} = 0
  \label{eq:nograd}
\end{equation}
zu \emph{jedem} Zeitpunkt. Der Vorfaktor hängt nur von $t$ ab, nicht vom Ort,
und fällt beim Betragsquadrat heraus.

Gl.~\eqref{eq:centering} ist hinreichend. Sie ist auch notwendig, solange die
Basisfunktionen linear unabhängig sind -- was für getrennte Spotpositionen
generisch der Fall ist; bei zusätzlichen Entartungen der Anordnung wären
weitere Lösungen denkbar.

\paragraph{Numerische Kontrolle.}
Geprüft wurde nicht nur die Folgerung, sondern Gl.~\eqref{eq:nograd} selbst.
Bei festem $t$ und über verschiedene Orte $r$ muss der Quotient
$E(-r,t)/E^{*}(r,t)$ betragsmäßig eins und in der Phase ortsunabhängig sein:

\begin{center}
\begin{tabular}{lcc}
  \toprule
  Phasensatz & $\bigl|E(-r,t)/E^{*}(r,t)\bigr|$ & Streuung von $\arg$ über $r$\\
  \midrule
  konjugiert, $c_x=c_y=90^\circ$ & $1{,}000000$ & $2{,}7\cdot10^{-6}\,\mathrm{rad}$\\
  gewählter Arbeitspunkt         & $1{,}000000$ & $2{,}5\cdot10^{-7}\,\mathrm{rad}$\\
  nur gespiegelt                 & $0{,}91 \pm 0{,}55$ & $7{,}5\,\mathrm{rad}$\\
  \bottomrule
\end{tabular}
\end{center}

Ebenso bestätigt sind die drei Voraussetzungen am Arbeitspunkt:
$x_n+x_{N-1-n}=0$ und $y_m+y_{M-1-m}=0$ bis auf
\SI{5e-8}{\micro\meter} Gitterauflösung, $A_n=[0{,}985;\,1;\,0{,}985]$ und
$B_m=[1{,}077;\,1;\,1;\,1{,}077]$ symmetrisch, sowie
$f_n+f_{N-1-n} = g_m+g_{M-1-m} = \SI{200.37}{\mega\hertz}$ konstant.

\subsection{Der entscheidende Unterschied: konjugiert, nicht gespiegelt}

\paragraph{Die beiden Begriffe.}
Für ein Spiegelpaar $(n,\bar n)$ stehen zwei Bedingungen zur Wahl:
\begin{align}
  \text{\emph{gespiegelt}:}&&  \varphi_{\bar n} &= \varphi_n
  &&\Longleftrightarrow&& e^{\,\mathrm{i}\varphi_{\bar n}} &= e^{\,\mathrm{i}\varphi_{n}} ,\\
  \text{\emph{konjugiert}:}&&  \varphi_{\bar n} &= c - \varphi_n
  &&\Longleftrightarrow&& e^{\,\mathrm{i}\varphi_{\bar n}} &= e^{\,\mathrm{i}c}\bigl(e^{\,\mathrm{i}\varphi_{n}}\bigr)^{*} .
\end{align}
Komplexe Konjugation spiegelt einen Zeiger in der Zahlenebene an der reellen
Achse, $\varphi\rightarrow-\varphi$; der Faktor $e^{\mathrm{i}c}$ dreht diese
Spiegelachse auf den Winkel $c/2$. Schreibt man
$\varphi_n = c/2 + \chi_n$, so lautet die Bedingung schlicht
\begin{equation}
  \chi_{\bar n} \;=\; -\,\chi_n ,
\end{equation}
die beiden Zeiger eines Spiegelpaares stehen also \emph{gleich weit vom
Mittenwinkel $c/2$ entfernt, nur nach entgegengesetzten Seiten}. Kurz:
\begin{center}
\begin{tabular}{@{}p{0.20\textwidth}p{0.72\textwidth}@{}}
  \toprule
  \textbf{gespiegelt} & gleicher Winkel\\
  \textbf{konjugiert} & gleicher \emph{Abstand} zu einem Bezugswinkel,
  entgegengesetztes Vorzeichen\\
  \bottomrule
\end{tabular}
\end{center}

Am gewählten Arbeitspunkt sieht das so aus:
\begin{align}
  \varphi_y &= [\,0;\;98;\;16;\;114\,]^\circ,
  & c_y &= 114^\circ,
  & c_y/2 &= 57^\circ, \nonumber\\
  \chi_y    &= [\,-57;\;+41;\;-41;\;+57\,]^\circ
  &&&& \text{antisymmetrisch.}
\end{align}
Gespiegelt ist dieser Satz gerade \emph{nicht} ($0^\circ \neq 114^\circ$).
Auf der $x$-Achse ist $\varphi_x=[\,0;\;12{,}5;\;25\,]^\circ$ mit
$c_x/2 = 12{,}5^\circ$ und $\chi_x=[\,-12{,}5;\;0;\;+12{,}5\,]^\circ$ -- bei
ungeradem $N$ sitzt der mittlere Ton notwendig auf dem Mittenwinkel, er ist
sein eigener Spiegelpartner.

\paragraph{Warum gerade die Konjugation.}
Der Grund steckt in Gl.~\eqref{eq:symfreq}. Die Frequenzspiegelung bedeutet
\begin{equation}
  e^{\,\mathrm{i}2\pi f_{\bar n}t}
  \;=\; e^{\,\mathrm{i}2\pi F t}\,\bigl(e^{\,\mathrm{i}2\pi f_{n}t}\bigr)^{*} ,
\end{equation}
die \emph{dynamische} Phase des Spiegelpartners konjugiert sich also von
selbst. Damit der gesamte Partnerterm das Konjugierte des Originals wird --
und nur dann kommt $|E(-r,t)|=|E(r,t)|$ heraus --, muss die \emph{statische}
Phase dasselbe tun. Genau das fordert Gl.~\eqref{eq:centering}.

\paragraph{Der Sonderfall, der beides erfüllt.}
Eine reelle Zahl ist gleich ihrem Konjugierten. Für
$\varphi\in\{0,\pi\}$ fallen beide Bedingungen deshalb zusammen -- daher die
Verwechslungsgefahr. Die naheliegende Vermutung, dass allgemein
\emph{spiegelsymmetrische} Phasen $\varphi_{\bar n} = \varphi_n$ genügen, ist
falsch, und der Unterschied ist messbar. Tabelle~\ref{tab:mirror} stellt beides gegenüber.

\begin{table}[htb]
  \centering
  \caption{Maximaler Schwerpunktversatz des Interferenzflecks über eine volle
  Grundperiode, $\sigma = \SI{108}{\nano\meter}$.}
  \label{tab:mirror}
  \begin{tabular}{llc}
    \toprule
    Phasensatz (Grad) & Eigenschaft & max.\ Versatz\\
    \midrule
    $\varphi_x=[0,90,0]$, $\varphi_y=[0,90,90,0]$      & gespiegelt             & \SI{43.65}{\nano\meter}\\
    $\varphi_x=[0,120,0]$, $\varphi_y=[0,50,50,0]$     & gespiegelt             & \SI{83.95}{\nano\meter}\\
    $\varphi_x=[0,45,90]$, $\varphi_y=[0,30,60,90]$    & konjugiert ($c=90$)    & \SI{0.00}{\nano\meter}\\
    $\varphi_x=[0,60,120]$, $\varphi_y=[0,140,-40,100]$& konjugiert ($120/100$) & \SI{0.00}{\nano\meter}\\
    $\varphi_x=[0,180,0]$, $\varphi_y=[0,180,180,0]$   & beides (reelle Zeiger) & \SI{0.00}{\nano\meter}\\
    \bottomrule
  \end{tabular}
\end{table}

Der letzte Fall erklärt, warum die Verwechslung naheliegt: Für
$\varphi\in\{0,\pi\}$ ist der Zeiger reell, Spiegelung und Konjugation fallen
zusammen, und beide Bedingungen sind gleichzeitig erfüllt. Sobald eine Phase
einen anderen Wert annimmt, trennen sich die Begriffe.

\paragraph{Anschaulich, am Spotpaar.}
Die Intensität ist die Summe der Kreuzterme aller Spotpaare,
\begin{equation}
  2\,A_s A_{s'}\, u(r-c_s)\,u(r-c_{s'})\,
  \cos\!\bigl(2\pi \Delta f\, t + \Delta\varphi\bigr),
  \qquad
  \Delta f = f_s - f_{s'},\;\; \Delta\varphi = \varphi_s-\varphi_{s'} .
\end{equation}
Der Ortsanteil ist um den Mittelpunkt $(c_s+c_{s'})/2$ zentriert, also im
Allgemeinen \emph{nicht} um den Ursprung. Er wird erst dadurch ausgeglichen,
dass es das gespiegelte Paar $(\bar s,\bar s')$ gibt, dessen Ortsanteil das
Spiegelbild ist. Die beiden Lappen heben sich aber nur dann gegenseitig zu
einem symmetrischen Muster auf, wenn sie \emph{im Gleichtakt} atmen. Für das
gespiegelte Paar ist $\Delta f \rightarrow -\Delta f$, und
\begin{itemize}\itemsep0.2em
  \item mit Gl.~\eqref{eq:centering}: $\Delta\varphi\rightarrow-\Delta\varphi$,
        also $\cos(-2\pi\Delta f t-\Delta\varphi)=\cos(2\pi\Delta f t+\Delta\varphi)$
        -- \textbf{identischer} Zeitverlauf, die Lappen steigen und fallen
        gemeinsam;
  \item bei bloßer Spiegelung: $\Delta\varphi\rightarrow+\Delta\varphi$, also
        $\cos(2\pi\Delta f t-\Delta\varphi)\neq\cos(2\pi\Delta f t+\Delta\varphi)$
        -- die beiden Lappen laufen \emph{gegeneinander}, und das Muster
        schwappt seitlich hin und her.
\end{itemize}
Nur für $\Delta\varphi\in\{0,\pi\}$ stimmen beide überein -- wieder der
Sonderfall reeller Zeiger.

\subsection{Struktur der Familie}

Mit $\varphi_n = c_x/2 + \chi_n$ geht Gl.~\eqref{eq:centering} über in
\begin{equation}
  \chi_n \;=\; -\,\chi_{N-1-n},
  \label{eq:antisym}
\end{equation}
die Phasen sind also eine globale Phase plus einen \emph{antisymmetrischen}
Anteil. Jeder Ton ist damit genauso weit von der Mittenphase $c_x/2$ entfernt
wie sein Spiegelpartner, nur in die andere Richtung -- das Gegenstück zur
Frequenzspiegelung aus Gl.~\eqref{eq:symfreq}. Die antisymmetrischen Vektoren
im $\mathbb{R}^N$ bilden einen Unterraum der Dimension $\lfloor N/2\rfloor$;
nach Abzug der global irrelevanten Phase bleiben
\begin{equation}
  \dim(\text{Familie}) = \Bigl\lfloor \tfrac{N}{2} \Bigr\rfloor
  \quad\text{statt}\quad
  N-1 .
\end{equation}
\textbf{Das Kriterium halbiert die Phasenfreiheit.} Bei ungeradem $N$
scheint eine diskrete Wahl hinzuzukommen, weil
$\chi_{(N-1)/2}=-\chi_{(N-1)/2}$ modulo $2\pi$ die beiden Lösungen $0$ und
$\pi$ zulässt. Sie ist jedoch keine zusätzliche Dimension: Verschiebt man die
Rampensteigung um $2\pi/(N-1)$, so bleibt $c$ modulo $2\pi$ unverändert,
während sich die mittlere Phase um genau $\pi$ ändert. Die beiden Lösungen
liegen also bereits auf demselben stetigen Ast
(Abschnitt~\ref{sec:crestback}).

Numerisch geprüft: zwölf zufällige Sätze aus dieser Familie ergeben über die
gesamte Grundperiode einen Schwerpunktversatz unter
$1{,}2\cdot10^{-4}\,\mathrm{nm}$; wird Gl.~\eqref{eq:centering} auf nur einer
Achse verletzt, sind es \SIrange{16}{74}{\nano\meter}.

\subsection{Der mittlere Ton ist sein eigener Spiegelpartner}
\label{sec:selfpair}

Gegen Gl.~\eqref{eq:centering} lässt sich ein naheliegender Einwand
erheben, und er ist wichtig genug für einen eigenen Abschnitt:

\begin{center}
\emph{„Der mittlere Ton hat doch gar keinen Partner -- also gibt es für
seine Phase auch keine Bedingung.``}
\end{center}

\noindent Der Einwand ist falsch, aber aus einem Grund, den man leicht
übersieht. Die Spiegelung wirkt auf den \emph{Indizes}:
\begin{equation}
  n \;\longmapsto\; \bar n = N-1-n .
\end{equation}
Das ist eine Involution ($\bar{\bar n} = n$) auf der Menge
$\{0,\dots,N-1\}$, und Involutionen haben nicht nur Paare, sondern
möglicherweise auch \emph{Fixpunkte}:

\begin{center}
\begin{tabular}{cll}
  \toprule
  $N$ & Zuordnung $n\mapsto\bar n$ & Fixpunkt\\
  \midrule
  3 & $0\!\leftrightarrow\!2$, \; $1\!\mapsto\!1$ & $n=1$\\
  4 & $0\!\leftrightarrow\!3$, \; $1\!\leftrightarrow\!2$ & keiner\\
  5 & $0\!\leftrightarrow\!4$, \; $1\!\leftrightarrow\!3$, \; $2\!\mapsto\!2$ & $n=2$\\
  \bottomrule
\end{tabular}
\end{center}

\noindent Ein Fixpunkt ist \emph{kein partnerloses} Element -- er ist ein
Element, das mit \emph{sich selbst} gepaart ist. Und Gl.~\eqref{eq:centering}
verlangt die Bedingung für \emph{jedes} $n$, den Fixpunkt eingeschlossen.
Dort liest sie sich
\begin{equation}
  \varphi_{\bar n} + \varphi_n \;=\; \varphi_n + \varphi_n \;=\; 2\varphi_n \;=\; c
  \qquad\text{für } n = \tfrac{N-1}{2} ,
\end{equation}
und damit ist die mittlere Phase festgelegt:
\begin{empheq}[box=\fbox]{equation}
  \varphi_{(N-1)/2} = \frac{c}{2} \quad\text{oder}\quad \frac{c}{2}+\pi
  \qquad (N \text{ ungerade}).
  \label{eq:midphase}
\end{empheq}
Die zwei Lösungen kommen daher, dass Gl.~\eqref{eq:centering} modulo $2\pi$
gilt; $2\varphi = c$ hat dann eben zwei Wurzeln. Beide gehören zur Familie --
und sie sind nicht so getrennt, wie sie aussehen: Bei festem $c$ liegen sie
$\pi$ auseinander, aber $c$ selbst ist ein Parameter, und ein Weg durch die
Familie führt stetig von der einen Wurzel zur anderen
(Abschnitt~\ref{sec:crestback}).

\paragraph{Warum die Bedingung nicht entfallen kann.}
Der Kern des Einwands ist die stillschweigende Annahme, jedes Paar dürfe
seine eigene Konstante haben. Genau das ist nicht der Fall. Zurück zum
Ursprung: Verlangt war
\begin{equation}
  E(-\mathbf{r},t) = e^{\mathrm{i}\alpha(t)}\,E^{*}(\mathbf{r},t) ,
\end{equation}
mit \emph{einem einzigen} Faktor $e^{\mathrm{i}\alpha(t)}$ für das
\emph{ganze} Feld. In der Herleitung wird der Beitrag des Spots $s$ auf der
linken Seite gegen den Beitrag des Spiegelspots $\bar s$ auf der rechten
gehalten; daraus wird $\varphi_s + \varphi_{\bar s} = c$. Der Wert $c$ ist
dabei nichts anderes als $\alpha$ -- und $\alpha$ ist für alle Spots
dasselbe, sonst addierten sich die Beiträge gar nicht erst zu einem
konjugierten Feld.

Der mittlere Ton sitzt bei der Mittenfrequenz $\bar f = F/2$, sein Spot also
genau im Spiegelzentrum. Unter der Spiegelung geht er in sich selbst über,
und seine Bedingung lautet daher nicht „keine``, sondern „mit sich selbst``:
Sein Beitrag muss gleich seinem eigenen Konjugierten mal demselben
$e^{\mathrm{i}\alpha}$ sein. Das ist eine Gleichung, und sie bindet
$\varphi_{(N-1)/2}$ an dasselbe $c$ wie alle anderen. \emph{Hätte} jedes Paar
seine eigene Konstante, wäre der mittlere Ton tatsächlich frei -- dann wäre
aber auch die Symmetrie des Feldes dahin.

\paragraph{Gegenprobe am Rechner.}
Das lässt sich unmittelbar prüfen. Ist das Kriterium erfüllt, so ist der
Intensitätsgradient am Atomort zu \emph{jedem} Zeitpunkt null
(Gl.~\eqref{eq:nograd}). Mit $\varphi_x = (0,\psi,0)$ ist $c = 0$, das
Kriterium verlangt also $\psi = 0$ oder $180^\circ$. Die $y$-Phasen bleiben
dabei unverändert auf $(0;98;16;114)$, erfüllen ihr eigenes Kriterium
($0+114 = 98+16$) und werden nicht angetastet:

\begin{table}[htb]
  \centering
  \caption{Größter Intensitätsgradient am Atomort über eine volle
  Grundperiode, normiert auf die mittlere Intensität. Die Werte bei
  $\psi = 0$ und $180^\circ$ sind das Rauschen der numerischen Ableitung.}
  \label{tab:midphase}
  \begin{tabular}{crrl}
    \toprule
    $\psi$ & $\max_t|\partial_x I|/I$ & $\max_t|\partial_y I|/I$ & \\
    \midrule
    $0^\circ$   & $4\cdot10^{-10}$ & $4\cdot10^{-10}$ & Kriterium erfüllt\\
    $10^\circ$  & $8{,}8\cdot10^{-4}$ & $2{,}6\cdot10^{-4}$ & \\
    $30^\circ$  & $2{,}5\cdot10^{-3}$ & $7{,}4\cdot10^{-4}$ & \\
    $60^\circ$  & $4{,}4\cdot10^{-3}$ & $1{,}3\cdot10^{-3}$ & \\
    $90^\circ$  & $5{,}1\cdot10^{-3}$ & $1{,}5\cdot10^{-3}$ & \\
    $150^\circ$ & $2{,}5\cdot10^{-3}$ & $7{,}4\cdot10^{-4}$ & \\
    $180^\circ$ & $4\cdot10^{-10}$ & $4\cdot10^{-10}$ & Kriterium erfüllt\\
    \bottomrule
  \end{tabular}
\end{table}

\noindent Der Gradient verschwindet exakt bei den beiden von
Gl.~\eqref{eq:midphase} vorhergesagten Werten und bei keinem anderen. Mit
$c = 60^\circ$ statt $0$, also $\varphi_x = (0,\psi,60^\circ)$, wandern die
Nullstellen mit auf $\psi = 30^\circ$ und $210^\circ$ -- genau $c/2$ und
$c/2+\pi$. Von „keiner Bedingung`` kann also keine Rede sein; es ist eine
Bedingung mit zwei Lösungen.

\paragraph{Nebenbefund: die Achsen sind nicht unabhängig.}
Auffällig an Tabelle~\ref{tab:midphase} ist die dritte Spalte. Verletzt wird
das Kriterium nur in $x$ -- die $y$-Phasen bleiben unverändert --, und
trotzdem kippt auch der $y$-Gradient. Der Grund ist die Einhüllende: Mit
einem Gaußprofil faktorisiert das Feld in $X(x,t)\,Y(y,t)$, und dann bleibt
$\partial_y I|_0 = 0$ tatsächlich erhalten (nachgerechnet: $3\cdot10^{-10}$
in allen Zeilen). Die reale Airy-Einhüllende ist \emph{radial} und damit
nicht separabel; sie koppelt die Achsen. Das Kriterium muss deshalb auf
\emph{beiden} Achsen gelten -- eine Verletzung in $x$ lässt sich nicht durch
die $y$-Phasen heilen.

\paragraph{Und was bei geradem \boldmath$N$ anders ist.}
Für gerades $N$ gibt es keinen Fixpunkt: Jeder Index hat einen von ihm
verschiedenen Partner, und keine einzelne Phase wird gebunden. Das ist der
strukturelle Grund für den Unterschied, der sich durch den ganzen Text zieht
-- bei $N_x = 3$ ist die mittlere Phase durch Gl.~\eqref{eq:midphase}
vergeben und der Crestfaktor damit festgenagelt
(Abschnitt~\ref{sec:psi}), während bei $N_y = 4$ ein freier Parameter bleibt,
mit dem sich $\crest_y$ von $2{,}83$ auf $2{,}20$ drücken lässt.

\subsection{Was die Symmetrie leistet}

Entwickelt man die Pulsfläche um den Atomort,
\begin{equation}
  \theta(r) \;\approx\; \theta(0) \;+\; \nabla\theta\big|_0\cdot r
  \;+\; \tfrac{1}{2}\,r^{\mathsf{T}} H\, r ,
  \label{eq:taylor}
\end{equation}
so liefert der lineare Term über eine Wolke der Breite $\sigma$ den Beitrag
$\Uw \approx |\nabla\theta|\,\sigma/\theta(0)$, der quadratische erst einen
Beitrag der Ordnung $\sigma^{2}$. Gl.~\eqref{eq:nograd} eliminiert den
linearen Term \emph{exakt und zu allen Zeiten}. Die Symmetrie beseitigt damit
genau den dominanten Beitrag zur Inhomogenität -- sie ist das eigentliche
Werkzeug, nicht ein Nebenaspekt.

Das lässt sich direkt nachweisen, indem man $\Uw$ als Funktion der Atomgröße
aufträgt: Ein linearer Gradient ergibt $\Uw\propto\sigma$, reine Krümmung
$\Uw\propto\sigma^{2}$.

\begin{table}[htb]
  \centering
  \caption{$\Uw$ gegen die Atomgröße. Der symmetrische Satz zeigt reine
  Krümmungsskalierung, der unsymmetrische Gradientenskalierung.}
  \label{tab:scaling}
  \begin{tabular}{ccc}
    \toprule
    $\sigma$ [\si{\nano\meter}] & symmetrisch (gewählter Punkt) & Schroeder\\
    \midrule
     27 & \SI{0.0123}{\percent} & \SI{3.55}{\percent}\\
     54 & \SI{0.0495}{\percent} & \SI{9.14}{\percent}\\
    108 & \SI{0.2063}{\percent} & \SI{17.92}{\percent}\\
    216 & \SI{1.0073}{\percent} & \SI{32.98}{\percent}\\
    \midrule
    Exponent $p$ in $\Uw\propto\sigma^{p}$ & $\mathbf{2{,}11}$ & $\mathbf{1{,}06}$\\
    \bottomrule
  \end{tabular}
\end{table}

Der gemessene Exponent bestätigt Gl.~\eqref{eq:nograd} unmittelbar: $p=1{,}06$
ohne Symmetrie, $p=2{,}11$ mit. Bei der tatsächlichen Atomgröße von
\SI{108}{\nano\meter} entspricht das einem Faktor $87$ in der Homogenität.

\subsection{Was die Symmetrie \emph{nicht} leistet: ein helles Zentrum}

Gl.~\eqref{eq:nograd} sagt, dass der Atomort ein \emph{stationärer} Punkt von
$\theta(r)$ ist -- nicht, dass er ein Maximum ist. Symmetrie schließt den
linearen Gradienten aus, unterscheidet aber nicht zwischen Maximum, Minimum
und Sattelpunkt. Tatsächlich durchläuft das Zentrum im Lauf einer
Grundperiode alle drei Fälle, wie Tabelle~\ref{tab:hell} am gewählten
Phasensatz zeigt.

\begin{table}[htb]
  \centering
  \caption{Das Zentrum bleibt stationär, wechselt aber im Lauf der Periode
  zwischen hell und dunkel. Ausgewählte Trigger-Zeitpunkte des gewählten
  Phasensatzes.}
  \label{tab:hell}
  \begin{tabular}{cccl}
    \toprule
    $t_0$ [\si{\micro\second}] & $A(t_0)/\langle A\rangle$ & $\Uw$ [\si{\percent}] & Zentrum ist\\
    \midrule
    $1{,}622$  & $2{,}69$ & $1{,}30$  & Sattelpunkt\\
    $9{,}730$  & $0{,}74$ & $0{,}62$  & Maximum\\
    $12{,}973$ & $1{,}17$ & $0{,}26$  & Maximum\\
    $6{,}486$  & $0{,}15$ & $48{,}99$ & \textbf{Minimum}\\
    $14{,}595$ & $0{,}09$ & $42{,}38$ & \textbf{Minimum}\\
    \bottomrule
  \end{tabular}
\end{table}

Über die gesamte Periode schwankt $A(t_0)$ zwischen $0{,}067$ und
$2{,}808\,\langle A\rangle$. Sitzt der Trigger auf einem Minimum, so ist der
Gradient zwar weiterhin exakt null, das Atom bekommt aber kaum Licht, und
$\Uw$ steigt auf knapp \SI{50}{\percent} -- nicht wegen eines Gradienten,
sondern wegen der \emph{Krümmung}: An einem schmalen Minimum wächst $\theta$
nach allen Seiten steil an, und über $\pm\sigma$ ist die Streuung
entsprechend groß. An einem breiten Maximum ist die Krümmung klein und $\Uw$
bleibt bei einigen Zehntelprozent. Im Mittel über die Periode beträgt $\Uw$
bei $A\ge\langle A\rangle$ nur \SI{2.4}{\percent}, andernfalls
\SI{11.1}{\percent}.

Daraus ergibt sich die Arbeitsteilung zwischen den beiden Parametergruppen:
\begin{center}
\begin{tabular}{@{}p{0.22\textwidth}p{0.70\textwidth}@{}}
  \toprule
  \textbf{Tonphasen} & legen die Symmetrie fest, also kein Gradient -- und
  das zu allen Zeiten\\
  \textbf{Trigger $t_0$} & wählt aus, \emph{welcher} stationäre Punkt: helles
  flaches Maximum statt dunkles Minimum\\
  \bottomrule
\end{tabular}
\end{center}
Deshalb führt die Optimierung die Nebenbedingung $A(t_0)\ge\langle A\rangle$
mit: Sie schließt die dunklen Extrema aus, die die Symmetrie allein nicht
verhindert.

\subsection{Rückwirkung auf den Crestfaktor}\label{sec:crestback}

Mit Gl.~\eqref{eq:crestsimple} lässt sich jetzt genau sagen, was die
Zentrierungsbedingung der HF-Kette antut. Der Nenner ist ohnehin
phasenunabhängig; zu klären ist allein, wie sich $\max_t|A|$ innerhalb der
Familie verhält. Für die $x$-Achse des Arbeitspunkts, $N=3$, geht das in
fünf Zeilen -- und das Ergebnis ist die zentrale Aussage dieses Abschnitts,
deshalb steht es zuerst.

\paragraph{\boldmath$N=3$: der Crestfaktor ist immer der schlechtestmögliche.}
Die Einhüllende dreier Töne ist nach
Gl.~\eqref{eq:envdef}, mit $\theta := 2\pi d\,t$ und den
\emph{Spannungs}amplituden
\begin{equation}
  a_0 = a_2 = \varrho := \sqrt{r_x} , \qquad a_1 = 1
  \qquad\text{(Konvention: } r_x \text{ ist ein \emph{Leistungs}verhältnis),}
  \label{eq:rhodef}
\end{equation}
\begin{equation}
  A(t) = \varrho\,e^{\mathrm{i}(\varphi_0-\theta)} + e^{\mathrm{i}\varphi_1}
       + \varrho\,e^{\mathrm{i}(\varphi_2+\theta)} .
\end{equation}
Das Zentrierungskriterium verlangt $\varphi_0+\varphi_2 = c$ und
$2\varphi_1 = c$ (Abschnitt~\ref{sec:selfpair}). Die erste Bedingung heißt,
dass die beiden Außentöne symmetrisch um $c/2$ liegen; man schreibt sie als
\begin{equation}
  \varphi_0 = \tfrac{c}{2}-\chi , \qquad
  \varphi_2 = \tfrac{c}{2}+\chi , \qquad
  \varphi_1 = \tfrac{c}{2} \;\;\text{oder}\;\; \tfrac{c}{2}+\pi ,
\end{equation}
mit $\chi$ als einzigem weiteren Parameter. Einsetzen und $e^{\mathrm{i}c/2}$
ausklammern:
\begin{align}
  A(t)\,e^{-\mathrm{i}c/2}
  &= \varrho\,e^{-\mathrm{i}(\chi+\theta)} \;\pm\; 1 \;+\; \varrho\,e^{+\mathrm{i}(\chi+\theta)}\\
  &= 2\varrho\cos(\theta+\chi) \;\pm\; 1 .
  \label{eq:n3direct}
\end{align}
Das Vorzeichen unterscheidet die beiden Zweige für $\varphi_1$. Und jetzt
steht alles da: $|A| = |\,2\varrho\cos(\theta+\chi)\pm1\,|$ hängt von $t$ nur noch
über den \emph{reellen} Kosinus ab, und der durchläuft über eine volle
Periode das ganze Intervall $[-1,1]$ -- unabhängig davon, was $\chi$ ist
($\chi$ legt nur fest, \emph{wann}). Das Maximum ist damit eine
Maximierung über $u := \cos(\theta+\chi) \in [-1,1]$:
\begin{align}
  \text{Zweig }+:&\quad \max_{u\in[-1,1]} |\,2\varrho\,u + 1\,| = |\,2\varrho+1\,| = 2\varrho+1
    &&\text{bei } u=+1 ,\\
  \text{Zweig }-:&\quad \max_{u\in[-1,1]} |\,2\varrho\,u - 1\,| = |\!-\!2\varrho-1\,| = 2\varrho+1
    &&\text{bei } u=-1 .
\end{align}
Beide Zweige liefern denselben Wert; der Minus-Zweig erreicht ihn nur am
anderen Ende des Intervalls. Mit $a_0=a_2=\varrho$ und $a_1=1$ ist
$2\varrho+1 = a_0+a_1+a_2$, also
\begin{equation}
  \max_t |A| = 2\varrho+1 = \sum_n a_n
  \qquad\text{für \emph{beide} Zweige und \emph{jedes} } \chi, c .
\end{equation}
Nach der Dreiecksungleichung ist $|A(t)|\le\sum_n a_n$ aber für
\emph{jeden überhaupt möglichen} Phasensatz die obere Schranke. Die
zentrierte Familie nimmt sie also stets an, und mit
Gl.~\eqref{eq:crestsimple} folgt
\begin{equation}
  \crest_x = \sqrt2\,\frac{\sum_n a_n}{\sqrt{\sum_n a_n^2}} .
  \label{eq:crestn3}
\end{equation}

\medskip
\noindent\textbf{Woher die Zahlen kommen.} Der Parameter $r_x = 0{,}97$ des
Arbeitspunkts ist ein \emph{Leistungs}verhältnis -- er gibt an, wie viel
Beugungsleistung die Randtöne gegenüber den inneren bekommen. Die
HF-\emph{Spannung} trägt davon die Wurzel (Abschnitt~\ref{sec:abgrenzung}),
und in die Einhüllende gehen Spannungen ein:
\begin{equation}
  \varrho = \sqrt{r_x} = \sqrt{0{,}97} = 0{,}9848858 ,
  \qquad a = [\,\varrho,\;1,\;\varrho\,] .
\end{equation}
Damit
\begin{align}
  \sum_n a_n &= 2\varrho+1 = 2{,}9697716 ,\\
  \sum_n a_n^2 &= 2\varrho^2+1 = 2\,r_x + 1 = 2{,}94
    \qquad\Longrightarrow\qquad
  \sqrt{\textstyle\sum_n a_n^2} = 1{,}7146428 ,
\end{align}
und eingesetzt in Gl.~\eqref{eq:crestn3}
\begin{equation}
  \crest_x = \sqrt2\;\frac{2{,}9697716}{1{,}7146428} = 2{,}4494263 .
\end{equation}
Bemerkenswert ist die Quadratsumme: $\sum a_n^2 = 2r_x+1$ enthält das
Leistungsverhältnis \emph{linear}, die Wurzel hebt sich weg -- der
Effektivwert ist also direkt die Wurzel aus der Gesamtleistung, wie es sein
muss.

\medskip
\noindent\textbf{Und nein, das ist nicht \boldmath$\sqrt6$.} Bei exakt
gleichen Amplituden wäre $\crest_x = \sqrt{2N} = \sqrt6 = 2{,}4494897$, der
Höchstwert aus Gl.~\eqref{eq:sqrt2n}. Der Arbeitspunkt liegt mit
$2{,}4494263$ um $6\cdot10^{-5}$ darunter, also \SI{0.003}{\percent} -- die
leichte Randabsenkung $r_x=0{,}97$ senkt Zähler und Nenner fast im gleichen
Maß. Auf vier Stellen gerundet sind beide $2{,}4494$; dass sie verschieden
sind, sieht man erst in der fünften. Der Unterschied ist physikalisch
bedeutungslos, aber er erklärt, warum hier $2{,}4494$ und nicht $\sqrt6$
steht.

\medskip
\noindent\fbox{\parbox{0.96\textwidth}{\textbf{Bei drei Tönen ist der
Crestfaktor unter der Zentrierungsbedingung immer der schlechtestmögliche --
unabhängig von $c$, von $\chi$ und vom Zweig. Diese Achse lässt sich durch
keine Phasenwahl entlasten.} Der Grund steht in Gl.~\eqref{eq:n3direct}: Der
einzige freie Parameter $\chi$ taucht nur \emph{innerhalb} des Kosinus auf.
Er kann verschieben, wann die Zeiger sich ausrichten, aber nicht
\emph{ob}.}}

\medskip
\noindent Damit ist die Frage für die $x$-Achse beantwortet, ohne jede
weitere Maschinerie. Der Rest dieses Abschnitts klärt, warum es bei $N=4$
anders ist und was der allgemeine Grund dahinter ist -- nämlich dass der
Parameter $\chi$ nichts anderes als eine \emph{Zeitverschiebung} ist und
solche Verschiebungen den Crestfaktor grundsätzlich nicht berühren.

\paragraph{\boldmath$N=4$: eine echte Freiheit bleibt.}

Hier liefert Gl.~\eqref{eq:centering} $\varphi_y=[\,0,\,b,\,c-b,\,c\,]$ mit
\emph{zwei} Parametern. Die Rampe ist der Sonderfall $b=c/3$; jedes andere
$b$ ist eine antisymmetrische Abweichung, die \emph{keine} Rampe ist und die
Einhüllende deshalb wirklich umformt. Die beiden inneren Töne lassen sich
dann nicht mehr gleichzeitig mit den äußeren ausrichten, $\max|A|$ sinkt unter
$\sum a_n$, und mit ihm der Crestfaktor. Bei $c=114^\circ$:

\begin{center}
\begin{tabular}{cccl}
  \toprule
  $b$ & $\varphi_y$ [Grad] & $\max|A|$ & $\crest_y$\\
  \midrule
  $38^\circ$ & $[0,\,38,\,76,\,114]$ & $4{,}154$ & $2{,}8265$ \quad(Rampe, $=\sum a_n$)\\
  $60^\circ$ & $[0,\,60,\,54,\,114]$ & $4{,}022$ & $2{,}7364$\\
  $\mathbf{98^\circ}$ & $[0,\,98,\,16,\,114]$ & $\mathbf{3{,}228}$ & $\mathbf{2{,}1963}$\\
  $114^\circ$ & $[0,\,114,\,0,\,114]$ & $3{,}639$ & $2{,}4761$\\
  \bottomrule
\end{tabular}
\end{center}

\noindent$\sum a_n = 4{,}154$ und $\sqrt{\sum a_n^2}=2{,}078$; der
Minimalwert $2{,}196$ wird bei $b=98^\circ$ erreicht.

\paragraph{Und was das für die Leistung heißt.}
Über Gl.~\eqref{eq:pavglimit} übersetzt sich der Crestfaktor direkt in
erlaubte mittlere HF-Leistung. Mit $P_{1\mathrm{dB}}=\SI{3.98}{\watt}$:
\begin{align}
  \crest_y = 2{,}826 \;&\rightarrow\; P_\mathrm{mittel,y} = 3{,}98/7{,}99 = \SI{0.50}{\watt},\\
  \crest_y = 2{,}485 \;&\rightarrow\; P_\mathrm{mittel,y} = 3{,}98/6{,}18 = \SI{0.65}{\watt},\\
  \crest_y = 2{,}196 \;&\rightarrow\; P_\mathrm{mittel,y} = 3{,}98/4{,}82 = \SI{0.83}{\watt}.
\end{align}
Die Wahl von $b$ auf der $y$-Achse ist damit \SI{65}{\percent} mehr
HF-Leistung wert -- und weil die $x$-Achse ohnehin festliegt, ist sie die
\emph{einzige} Stellschraube, die die Zentrierung übrig lässt.

\subsection{Der allgemeine Grund: Rampen sind Zeitverschiebungen}
\label{sec:rampen}

Die Rechnung für $N=3$ in Abschnitt~\ref{sec:crestback} kam ohne weitere
Hilfsmittel aus: Der freie Parameter $\chi$ stand dort einfach im Kosinus und
konnte dessen Ausschlag nicht ändern. Für allgemeines $N$ braucht es die
dahinterliegende Aussage, und sie erklärt zugleich die Spalte
„crest-wirksam`` in Tabelle~\ref{tab:dims}:

\begin{center}
\emph{Eine linear mit dem Tonindex wachsende Phase -- eine \textbf{Rampe} --
verschiebt die Einhüllende nur in der Zeit und ändert den Crestfaktor
überhaupt nicht.}
\end{center}

\noindent Von den $\lfloor N/2\rfloor$ Dimensionen der zentrierten Familie
ist damit eine von vornherein wirkungslos; es bleiben
$\lfloor N/2\rfloor - 1$ crest-wirksame. Für $N=3$ ist das null -- genau das
Ergebnis von oben, jetzt als Spezialfall. Hier die Begründung Schritt für
Schritt.

\medskip
\noindent\textbf{Schritt 1: die Einhüllende im mitgeführten Bezugssystem.}
Nach Gl.~\eqref{eq:envdef} ist
\begin{equation}
  A(t) = \sum_n a_n\,e^{\,\mathrm{i}[\,2\pi (f_n-\bar f)\,t + \varphi_n]} ,
  \qquad \bar f = \frac{1}{N}\sum_n f_n ,
\end{equation}
wobei die Mittenfrequenz $\bar f$ abgezogen ist, weil der Träger für den
Betrag $|A|$ keine Rolle spielt.

\medskip
\noindent\textbf{Schritt 2: das Frequenzraster.} Mit dem äquidistanten Raster
$f_n = f_\mathrm{off} + n\,d$ ($d$ = Tonabstand) ist
$\bar f = f_\mathrm{off} + d\,\tfrac{N-1}{2}$ und daher
\begin{equation}
  f_n - \bar f = d\Bigl(n-\tfrac{N-1}{2}\Bigr) \;\equiv\; d\,\hat n ,
  \qquad \hat n := n-\tfrac{N-1}{2} .
  \label{eq:nhat}
\end{equation}
Der \emph{zentrierte} Index $\hat n$ läuft symmetrisch um null:
$\hat n = -1,0,+1$ bei $N=3$ und $-\tfrac32,-\tfrac12,+\tfrac12,+\tfrac32$
bei $N=4$.

\medskip
\noindent\textbf{Schritt 3: die Rampe in derselben Zählung.} Eine Rampe ist
eine Phase, die linear mit dem Tonindex wächst:
\begin{equation}
  \varphi_n = \varphi_\mathrm{s} + \delta\,n ,
  \qquad\text{also}\qquad
  \varphi = [\,\varphi_\mathrm{s},\;\varphi_\mathrm{s}+\delta,\;
               \varphi_\mathrm{s}+2\delta,\;\dots\,] .
\end{equation}
$\delta$ ist der Zuwachs von Ton zu Ton, $\varphi_\mathrm{s}$ der Startwert.
Letzterer ist eine allen Tönen gemeinsame Phase; er zieht sich als Faktor
$e^{\mathrm{i}\varphi_\mathrm{s}}$ vor die Summe und ändert $|A|$ nicht. Auch
sonst ist nur wichtig, wie die Phasen \emph{relativ} zueinander liegen. Es
ist deshalb zweckmäßig, den Mittelwert
\begin{equation}
  \bar\varphi = \frac{1}{N}\sum_n \varphi_n
              = \varphi_\mathrm{s} + \delta\,\frac{N-1}{2}
\end{equation}
abzuspalten. Man definiert die mittelwertfreie Phase
\begin{empheq}[box=\fbox]{equation}
  \chi_n \;:=\; \varphi_n - \bar\varphi
         \;=\; \delta\Bigl(n-\frac{N-1}{2}\Bigr)
         \;=\; \delta\,\hat n .
  \label{eq:rampchi}
\end{empheq}
\textbf{Das ist die ganze Einführung von $\chi_n$}: die Rampe, um ihren
eigenen Mittelwert verschoben, damit sie in derselben zentrierten Zählung
$\hat n$ steht wie die Frequenzen in Gl.~\eqref{eq:nhat}. Ausgeschrieben:
\begin{align}
  N=3:&\quad \chi = \delta\,[\,-1,\;0,\;+1\,] , \\
  N=4:&\quad \chi = \delta\,[\,-\tfrac32,\;-\tfrac12,\;+\tfrac12,\;+\tfrac32\,] .
\end{align}

\medskip
\noindent\textbf{Und warum man $\varphi_n$ durch $\chi_n$ ersetzen darf.}
Streng genommen darf man das \emph{nicht} -- man klammert aus. Mit
$\varphi_n = \bar\varphi + \chi_n$ ist
\begin{equation}
  A(t) = \sum_n a_n\,e^{\,\mathrm{i}(2\pi d\,\hat n\,t + \bar\varphi + \chi_n)}
       = e^{\,\mathrm{i}\bar\varphi}\underbrace{
         \sum_n a_n\,e^{\,\mathrm{i}(2\pi d\,\hat n\,t + \chi_n)}}_{
         =:\;\tilde A(t)} ,
  \label{eq:factorglobal}
\end{equation}
denn $\bar\varphi$ hängt \emph{nicht} von $n$ ab und darf deshalb vor die
Summe. Das Ersetzen ist also kein Weglassen, sondern der Übergang zu einer
neuen Größe $\tilde A$, und es gilt
\begin{equation}
  A(t) = e^{\,\mathrm{i}\bar\varphi}\,\tilde A(t)
  \quad\Longrightarrow\quad
  |A(t)| = |\tilde A(t)| \;\;\text{für alle } t
  \quad\Longrightarrow\quad
  \crest(A) = \crest(\tilde A) ,
\end{equation}
weil $|e^{\mathrm{i}\bar\varphi}| = 1$ ist und in Gl.~\eqref{eq:crestsimple}
nur Beträge vorkommen. $A$ und $\tilde A$ sind \emph{verschiedene} komplexe
Funktionen -- sie unterscheiden sich um eine konstante Drehung --, aber sie
haben denselben Betragsverlauf und damit denselben Crestfaktor. Nur darauf
kommt es hier an; im weiteren wird mit $\tilde A$ gerechnet.

\medskip
\noindent Nebenbei fällt damit ein Ergebnis ab, das für sich steht. Der
Spiegelindex $\bar n = N-1-n$ hat in zentrierter Zählung
\begin{equation}
  \hat{\bar n} = (N-1-n) - \frac{N-1}{2}
               = -\Bigl(n-\frac{N-1}{2}\Bigr) = -\hat n ,
\end{equation}
und daher ist $\chi_{\bar n} = \delta\,\hat{\bar n} = -\delta\,\hat n =
-\chi_n$. Das ist genau die Antisymmetrie, die Gl.~\eqref{eq:antisym} für die
zentrierte Familie verlangt: \textbf{Jede Rampe erfüllt das
Zentrierungskriterium automatisch.}

\medskip
\noindent\textbf{Schritt 4: gibt es überhaupt eine Verschiebung?} Hier ist
Sorgfalt nötig, sonst wird die Aussage zirkulär. Es wird \emph{nicht}
behauptet, man dürfe $\tau$ passend definieren -- gefragt ist, ob ein solches
$\tau$ überhaupt \emph{existiert}. Die Frage lautet also:

\begin{center}
\emph{Gibt es \textbf{eine} Zahl $\tau$ -- dieselbe für alle Zeiten und alle
Töne -- mit $\tilde A_\delta(t) = \tilde A_0(t+\tau)$ für \textbf{alle} $t$?}
\end{center}

\noindent Mit $\chi_n=\delta\hat n$ lautet die linke Seite
\begin{equation}
  \tilde A_\delta(t) = \sum_n \bigl(a_n\,e^{\,\mathrm{i}\delta\hat n}\bigr)\,
    e^{\,\mathrm{i}2\pi d\,\hat n\,t} ,
  \label{eq:adelta}
\end{equation}
und $\tilde A_0(t)=\sum_n a_n e^{\mathrm{i}2\pi d\hat n t}$ ist eine feste
Funktion, in der kein $\delta$ vorkommt. Einsetzen von $t+\tau$ gibt
\begin{equation}
  \tilde A_0(t+\tau) = \sum_n \bigl(a_n\,e^{\,\mathrm{i}2\pi d\hat n\tau}\bigr)\,
    e^{\,\mathrm{i}2\pi d\,\hat n\,t} .
  \label{eq:a0shift}
\end{equation}

\medskip
\noindent\textbf{Koeffizientenvergleich ist hier erlaubt}, und das ist nicht
selbstverständlich: Beide Seiten sind Entwicklungen nach denselben Funktionen
$e^{\mathrm{i}2\pi d\hat n t}$, und diese sind zu verschiedenen $\hat n$
\emph{linear unabhängig} (verschiedene Frequenzen). Zwei solche Summen
stimmen als Funktionen genau dann überein, wenn alle Koeffizienten
übereinstimmen. Also
\begin{equation}
  a_n\,e^{\,\mathrm{i}\delta\hat n} \;=\; a_n\,e^{\,\mathrm{i}2\pi d\hat n\tau}
  \qquad\Longleftrightarrow\qquad
  2\pi d\,\hat n\,\tau \;\equiv\; \delta\,\hat n \pmod{2\pi}
  \quad\text{für jedes } n .
  \label{eq:tausystem}
\end{equation}

\medskip
\noindent\textbf{Und hier steckt der Inhalt.} Gl.~\eqref{eq:tausystem} ist ein
Gleichungs\emph{system}: $N$ Bedingungen, \emph{eine} Unbekannte $\tau$. So
etwas ist im Normalfall unlösbar -- man kann nicht eine Zahl so wählen, dass
sie vier Gleichungen gleichzeitig erfüllt, es sei denn, die Gleichungen sind
nicht unabhängig. Genau das ist hier der Fall: Auf beiden Seiten steht
derselbe Faktor $\hat n$, er kürzt sich für $\hat n\neq0$ heraus, und für
$\hat n=0$ steht ohnehin $0=0$. Alle $N$ Bedingungen reduzieren sich auf
\emph{dieselbe}:
\begin{empheq}[box=\fbox]{equation}
  \tau = \frac{\delta}{2\pi d}
  \qquad\Longrightarrow\qquad
  \tilde A_\delta(t) = \tilde A_0(t+\tau) \quad\text{für alle } t .
  \label{eq:rampshift}
\end{empheq}
\emph{Die Aussage ist die Lösbarkeit, nicht der Wert von $\tau$.} Dass am Ende
eine Formel für $\tau$ dasteht, ist das Ergebnis und nicht die Voraussetzung.

\paragraph{Die Umkehrung: Zeitverschiebung \boldmath$\Longleftrightarrow$ Rampe.}
Am deutlichsten wird das, wenn man die Rechnung rückwärts liest. Sei
$\chi_n$ ein \emph{beliebiger} mittelwertfreier Phasensatz. Derselbe
Koeffizientenvergleich verlangt dann
\begin{equation}
  2\pi d\,\hat n\,\tau \;\equiv\; \chi_n \pmod{2\pi}
  \qquad\Longrightarrow\qquad
  \tau \;=\; \frac{\chi_n}{2\pi d\,\hat n} \quad\text{für jedes } \hat n\neq0 .
\end{equation}
Rechts steht eine Zahl, die im Allgemeinen \emph{von $n$ abhängt}. Ein
gemeinsames $\tau$ existiert genau dann, wenn der Quotient $\chi_n/\hat n$ für
alle Töne derselbe ist -- und das heißt
\begin{equation}
  \chi_n \equiv \delta\,\hat n \pmod{2\pi} ,
\end{equation}
also: $\chi$ ist eine Rampe. Damit ist gezeigt:

\medskip
\noindent\fbox{\parbox{0.96\textwidth}{Ein Phasensatz ist \emph{genau dann}
eine Zeitverschiebung der Einhüllenden, wenn er eine Rampe ist (modulo
$2\pi$). Die Bedingung an $\tau$ ist keine Definition, sondern eine
Lösbarkeitsbedingung, und sie charakterisiert die Rampen vollständig.}}

\medskip
\noindent Wie eng das ist, zeigt eine Abzählung: Bei $N=4$ ist der Raum der
Phasensätze (nach Abzug der globalen Phase) drei\-dimensional, die
Zeitverschiebungen bilden darin eine \emph{ein}dimensionale Kurve. Ein
zufälliger Phasensatz ist also so gut wie nie eine Verschiebung -- die
Rampen treffen diese dünne Menge exakt.

\medskip
\noindent Das „modulo $2\pi$`` ist dabei nicht kosmetisch. Der Satz
$\varphi_x = [0;\,180^\circ;\,0]$ etwa sieht nach keiner Rampe aus, ist aber
eine (mit $\delta=180^\circ$, siehe unten), und der Test findet ihn auch als
solche: Er fällt mit $5\cdot10^{-16}$ bei
$\tau=\SI{-2.7027}{\micro\second}$ durch, also bei genau einer halben
Einhüllendenperiode.

\medskip
\noindent\textbf{Wo das $\delta$ geblieben ist.} Auf der rechten Seite von
Gl.~\eqref{eq:rampshift} steht kein $\delta$ mehr, und das ist gerade der
Inhalt der Aussage: Die Rampe ist nicht verschwunden, sie ist \emph{in das
Argument gewandert}. $\tilde A_0$ ist eine einzige, ein für alle Mal
festliegende Funktion; die ganze Rampenfamilie besteht aus nichts als ihren
Zeitverschiebungen, und $\delta$ steckt nur noch in der Angabe, um wie viel
verschoben wird. Zurückübersetzt auf die tatsächlichen Phasen bedeutet das
\begin{equation}
  |A_\delta(t)| = |A_0(t+\tau)| \qquad\text{für alle } t ,
\end{equation}
denn die globalen Faktoren $e^{\mathrm{i}\bar\varphi}$ auf beiden Seiten haben
Betrag eins.


\medskip
\noindent\textbf{Richtung und Zahlenprobe.} Das Vorzeichen: Was $\tilde A_0$
zur Zeit $t+\tau$ tut, tut $\tilde A_\delta$ bereits zur Zeit $t$ -- eine positive
Rampe lässt die Einhüllende also \emph{früher} eintreffen. Das Maximum, das
bei $\delta=0$ auf $t=0$ liegt, wandert nach $t=-\tau$. Nachgerechnet auf
$200\,001$ Zeitpunkten, komplexwertig, mit
$\max_t|\tilde A_\delta(t)-\tilde A_0(t+\tau)|$ relativ zur Signalspitze:

\begin{center}
\begin{tabular}{lrrl}
  \toprule
  Phasensatz & $\tau$ & Abweichung & Maximum bei\\
  \midrule
  Rampe $\delta=12{,}5^\circ$  & \SI{0.1877}{\micro\second} & $6\cdot10^{-16}$ & $t=\SI{-0.1877}{\micro\second}$\\
  Rampe $\delta=25^\circ$      & \SI{0.3754}{\micro\second} & $6\cdot10^{-16}$ & \\
  Rampe $\delta=90^\circ$      & \SI{1.3514}{\micro\second} & $6\cdot10^{-16}$ & $t=\SI{-1.3514}{\micro\second}$\\
  Rampe $\delta=175{,}6^\circ$ & \SI{2.6366}{\micro\second} & $1\cdot10^{-15}$ & \\
  Rampe $\delta=-40^\circ$     & \SI{-0.6006}{\micro\second}& $6\cdot10^{-16}$ & \\
  \midrule
  \emph{keine} Rampe, $\psi=30^\circ$ & \multicolumn{1}{c}{--} & $\mathbf{2{,}3\cdot10^{-1}}$ & existiert nicht\\
  \emph{keine} Rampe, $\psi=90^\circ$ & \multicolumn{1}{c}{--} & $\mathbf{6{,}1\cdot10^{-1}}$ & existiert nicht\\
  \bottomrule
\end{tabular}
\end{center}

\noindent Also Maschinengenauigkeit, nicht nur im Betrag, sondern in der
komplexen Amplitude; für $N=4$ mit $d=\SI{0.1233}{\mega\hertz}$ ebenso
($\le2\cdot10^{-15}$ bei $\delta = 38^\circ, 98^\circ, 250^\circ$). Die
Zeile $\delta=0$ ist bewusst nicht aufgeführt: Sie lautete
$\tilde A_0(t)=\tilde A_0(t+0)$ und prüfte gar nichts.

Die beiden unteren Zeilen zeigen, dass der Test überhaupt etwas
\emph{unterscheidet}. Dort wurde $\psi$ aus Abschnitt~\ref{sec:psi} auf den
mittleren Ton gelegt -- eine Phase, die nicht linear in $\hat n$ ist -- und
anschließend $\tau$ \emph{sowie} eine freie globale Phase numerisch so gewählt,
dass die Abweichung minimal wird. Selbst mit dieser Bestanpassung bleiben
\SI{23}{\percent} bzw.\ \SI{61}{\percent} der Signalspitze übrig: Es gibt
schlicht kein $\tau$, das die Einhüllende zur Deckung bringt. Eine Tabelle,
in der jede Zeile besteht, beweist wenig -- diese hier kann auch durchfallen.

\medskip
\noindent\textbf{Schritt 5: und das ändert den Crestfaktor wirklich nicht.}
Nach Gl.~\eqref{eq:crestsimple} ist
$\crest = \sqrt2\,\max_t|\tilde A|\,/\,\mathrm{rms}_t|\tilde A|$, und beide
Größen werden über eine \emph{volle} Periode $T$ gebildet. Die Substitution
$t' = t+\tau$ bildet ein volles Periodenintervall wieder auf ein volles ab,
also ist die \emph{Wertemenge}
\begin{equation}
  \bigl\{\,|\tilde A_\delta(t)| \;:\; t\in[0,T)\,\bigr\}
  \;=\;
  \bigl\{\,|\tilde A_0(t')| \;:\; t'\in[\tau,T+\tau)\,\bigr\}
  \;=\;
  \bigl\{\,|\tilde A_0(t')| \;:\; t'\in[0,T)\,\bigr\}
\end{equation}
dieselbe -- nur in anderer Reihenfolge durchlaufen. Maximum und quadratisches
Mittel hängen aber nur von der Wertemenge ab, nicht von der Reihenfolge.
Also $\crest(\delta)=\crest(0)$, und zwar exakt und nicht näherungsweise.

\medskip
\noindent\textbf{Eine Feinheit bei geradem \boldmath$N$.} Der Schritt setzt
voraus, dass $|\tilde A|$ die Periode $T=1/d$ hat. Für ungerades $N$ ist
$\hat n$ ganzzahlig und $\tilde A$ selbst periodisch. Für gerades $N$ ist
$\hat n$ halbzahlig, und dann ist
\begin{equation}
  \tilde A(t+T) = \sum_n a_n e^{\,\mathrm{i}(2\pi d\hat n t + 2\pi\hat n + \chi_n)}
                = -\,\tilde A(t) ,
\end{equation}
denn $e^{\mathrm{i}2\pi\hat n} = e^{\mathrm{i}\pi(2n-N+1)} = -1$ für gerades
$N$. $\tilde A$ ist dort also \emph{anti}periodisch mit $T$ und erst mit $2T$
periodisch. Der Betrag $|\tilde A|$ hat aber in beiden Fällen die Periode $T$,
und da in den Crestfaktor nur Beträge eingehen, gilt das Argument unverändert.
Numerisch: $\tilde A(t+T)/\tilde A(t) = +1$ bei $N=3$, $-1$ bei $N=4$, und
$|\tilde A|$ in beiden Fällen $T$-periodisch.

\medskip
\noindent\textbf{Zahlenprobe.} Über $400\,000$ Stützstellen einer vollen
Periode, für die Rampensteigungen $\delta = 0^\circ$ bis $359^\circ$:

\begin{center}
\begin{tabular}{lccc}
  \toprule
  & $\max|\tilde A|$ & $\mathrm{rms}|\tilde A|$ & $\crest$\\
  \midrule
  $x$-Achse, $N=3$, jedes $\delta$ & $2{,}96977156$ & $1{,}71464282$ & $2{,}44942630$\\
  $y$-Achse, $N=4$, jedes $\delta$ & $4{,}15406592$ & $2{,}07846097$ & $2{,}82648385$\\
  \bottomrule
\end{tabular}
\end{center}

\noindent Übereinstimmung auf elf Stellen, einschließlich $\delta=180^\circ$
und $270^\circ$. Der Effektivwert trifft dabei exakt
$\sqrt{\sum_n a_n^2}$ aus Abschnitt~\ref{sec:rmsphase} -- er ist ja ohnehin
phasenunabhängig; die ganze Aussage dieses Abschnitts betrifft allein den
Zähler $\max_t|\tilde A|$.

\paragraph{Was eine Rampe sehr wohl ändert.}
Crestneutral heißt nicht folgenlos. Die Verschiebung ist gegenüber dem
\emph{absoluten} Zeitnullpunkt real, und der Trigger $t_0$ ist genau so ein
Nullpunkt. Hält man ihn fest und misst die Pulsfläche
$\int_{t_0}^{t_0+\SI{1}{\micro\second}}|A|^2\,\mathrm{d}t$ bei $t_0=0$:

\begin{center}
\begin{tabular}{cc}
  \toprule
  $\delta$ & Pulsfläche relativ\\
  \midrule
  $0^\circ$      & $1{,}000$\\
  $12{,}5^\circ$ & $0{,}825$\\
  $90^\circ$     & $0{,}043$\\
  $180^\circ$    & $0{,}066$\\
  \bottomrule
\end{tabular}
\end{center}

\noindent Eine Rampe von $90^\circ$ je Ton lässt bei unverändertem Trigger
nur noch \SI{4}{\percent} der Pulsfläche übrig -- sie schiebt den Puls in ein
Minimum der Einhüllenden. \emph{Redundant zu $t_0$} heißt also: Die Rampe
muss \emph{gemeinsam} mit $t_0$ betrachtet werden, nicht dass sie folgenlos
wäre. Für die Optimierung ist sie deshalb kein zusätzlicher Freiheitsgrad,
sondern eine zweite Schreibweise für einen vorhandenen.

\paragraph{Und ist das Interferenzmuster damit auch nur zeitverschoben?}
Nur wenn beide Achsen \emph{passend} gerampt werden. Jede Achse hat ihren
eigenen Tonabstand, $d_x=\SI{0.185}{\mega\hertz}$ und
$d_y=\SI{0.1233}{\mega\hertz}$, und damit ihren eigenen Versatz
$\tau_{x,y}=\delta_{x,y}/(2\pi d_{x,y})$. Nur für
\begin{equation}
  \frac{\delta_x}{d_x} = \frac{\delta_y}{d_y}
  \qquad\Longleftrightarrow\qquad \tau_x = \tau_y
\end{equation}
wandern beide Einhüllenden um dasselbe $\tau$, und erst dann ist auch das
zweidimensionale Intensitätsmuster $I(\mathbf r,t)$ als Ganzes verschoben.
Nachgerechnet mit $\delta_x = 30^\circ$ und $\tau=\SI{0.4505}{\micro\second}$,
über eine volle Grundperiode und das ganze Ortsgitter,
$\max|I_\mathrm{Rampe}(\mathbf r,t)-I_0(\mathbf r,t-\tau)|/\max I$:
\begin{align}
  \text{angepasste Rampen auf beiden Achsen:}&\quad 8\cdot10^{-14} ,\\
  \text{Rampe nur auf der $x$-Achse:}&\quad 0{,}28 .
\end{align}
\textbf{„Rampe ist Zeitverschiebung`` ist also eine Aussage über die
Einhüllende \emph{einer} HF-Achse -- und nur als solche wird sie hier
gebraucht, denn der Crestfaktor ist ebenfalls eine Größe je Achse.} Für das
Licht am Atom gilt sie erst bei passenden Rampen auf beiden Achsen.

\medskip
\noindent\textbf{Schritt 6: wie weit die Verschiebung reicht.} Durchläuft
$\delta$ einmal $[0,2\pi)$, so durchläuft $\tau = \delta/(2\pi d)$ genau
$[0,1/d)$ -- und $1/d$ ist die Periode der Einhüllenden. \emph{Die
Rampenfamilie deckt also exakt eine Einhüllendenperiode ab, nicht mehr und
nicht weniger.} Damit ist sie vollständig redundant zum Trigger $t_0$: Jede
Rampe lässt sich durch ein anderes $t_0$ ersetzen und umgekehrt. Am
Arbeitspunkt ($d=\SI{0.185}{\mega\hertz}$) gehört zu $\delta=12{,}5^\circ$
der Versatz $\tau=\SI{0.188}{\micro\second}$, zu einer vollen Rampe
$\tau=\SI{5.405}{\micro\second}$.


\paragraph{Nachtrag: die \boldmath$N=3$-Familie \emph{ist} die Rampenfamilie.}
Die beiden Zweige aus Abschnitt~\ref{sec:selfpair},
$\varphi_x=[\,0,\,c/2,\,c\,]$ und $\varphi_x=[\,0,\,c/2+\pi,\,c\,]$, sind zusammen
genau die Rampenfamilie -- ein Punkt, an dem man sich leicht verrechnet, weil die
Phasen nur \emph{modulo} $2\pi$ definiert sind. Die Rampe
$\varphi_n = \delta\,n$ liefert
\begin{equation}
  \varphi_x = [\,0,\;\delta,\;2\delta\,],
  \qquad c = \varphi_0+\varphi_2 = 2\delta,
  \qquad \varphi_1 = \delta .
\end{equation}
Durchläuft $\delta$ einmal $[0,2\pi)$, so durchläuft $c = 2\delta \bmod 2\pi$
den Bereich $[0,2\pi)$ \emph{zweimal}: Für $\delta<\pi$ ist $\varphi_1 = c/2$
(erster Zweig), für $\delta\ge\pi$ ist $\varphi_1 = c/2+\pi$ (zweiter). Die
scheinbar diskrete Zusatzwahl ist also nichts weiter als die zweite Hälfte
des Rampenkreises:

\begin{center}
\begin{tabular}{crcl}
  \toprule
  $\delta$ & $\varphi_x \bmod 360^\circ$ & $c$ & Zweig\\
  \midrule
  $45^\circ$  & $[0;\;45;\;90]$  & $90^\circ$  & $\varphi_1 = c/2$\\
  $90^\circ$  & $[0;\;90;\;180]$ & $180^\circ$ & $\varphi_1 = c/2$\\
  $225^\circ$ & $[0;\;225;\;90]$ & $90^\circ$  & $\varphi_1 = c/2+180^\circ$\\
  $270^\circ$ & $[0;\;270;\;180]$& $180^\circ$ & $\varphi_1 = c/2+180^\circ$\\
  \bottomrule
\end{tabular}
\end{center}

\noindent Zentrierte Familie und Rampenfamilie fallen bei $N=3$ also
vollständig zusammen. Das ist keine neue Information -- der Crestfaktor stand
schon in Abschnitt~\ref{sec:crestback} fest --, aber es erklärt, warum dort
nichts mehr zu holen war: Die einzige Freiheit der Familie ist eine
Verschiebung, und Verschiebungen sieht der Crestfaktor nicht.

\subsection{Freiheitsgrade im Überblick}

\begin{table}[htb]
  \centering
  \caption{Freiheitsgrade und erreichbarer Crestfaktor. $\crest_{\min}$
  jeweils numerisch über die zulässige Menge minimiert, gleiche Amplituden.}
  \label{tab:dims}
  \begin{tabular}{ccccccc}
    \toprule
    $N$ & frei: $N-1$ & Familie: $\lfloor N/2\rfloor$ & crest-wirksam
        & $\crest_{\min}$ (Familie) & $\crest_{\min}$ (frei) & $\sqrt{2N}$\\
    \midrule
     3 &  2 & 1 & \textbf{0} & $\mathbf{2{,}4495}$ & $1{,}8257$ & $2{,}4495$\\
     4 &  3 & 2 & 1 & $2{,}1773$ & $1{,}7481$ & $2{,}8284$\\
     5 &  4 & 2 & 1 & $2{,}0555$ & $1{,}7311$ & $3{,}1623$\\
     6 &  5 & 3 & 2 & $2{,}0077$ & $1{,}7568$ & $3{,}4641$\\
     8 &  7 & 4 & 3 & $2{,}0452$ & $1{,}6200$ & $4{,}0000$\\
    13 & 12 & 6 & 5 & $2{,}0143$ & $1{,}7016$ & $5{,}0990$\\
    \bottomrule
  \end{tabular}
\end{table}

Tabelle~\ref{tab:dims} zeigt den entscheidenden Punkt. Bei $N=3$ ist
$\lfloor N/2\rfloor-1=0$: Es bleibt kein crestwirksamer Parameter übrig, und
$\crest_x = \sqrt{6} = 2{,}4494$ ist Minimum \emph{und} Maximum der Familie
zugleich -- exakt der Höchstwert nach Gl.~\eqref{eq:sqrt2n}. Das ist
dasselbe Ergebnis wie in Abschnitt~\ref{sec:crestback}, hier nur aus der
Abzählung statt aus der expliziten Rechnung.

Ab $N=4$ öffnet sich die Schere: Während $\sqrt{2N}$ weiter wächst, fällt
$\crest_{\min}$ innerhalb der Familie auf etwa $2{,}0$ und bleibt dort. Der
\emph{Preis} des Kriteriums, also der Abstand zum unbeschränkten Optimum bei
rund $1{,}7$, schrumpft von „alles`` bei $N=3$ auf etwa $0{,}3$ ab $N=5$. Mit
mehr Tönen wird die exakte Zentrierung also nahezu gratis.

Mit den tatsächlichen Amplituden des Arbeitspunkts ($r_x=0{,}97$,
$r_y=1{,}16$) ergibt sich dasselbe Bild: $\crest_x$ ist auf $2{,}4494$
festgenagelt, $\crest_y$ lässt sich innerhalb der Familie von $2{,}8265$ auf
$2{,}1963$ absenken, unbeschränkt wären $1{,}7661$ erreichbar.

\subsection{Die eine freie Phase bei \boldmath$N=3$ -- und was sie kostet}
\label{sec:psi}

Der vorige Abschnitt endet mit einem Befund, der zunächst wie ein
Rechenfehler aussieht: Bei drei Tönen ist die mittlere Phase $\varphi_1$ doch
eine ganz gewöhnliche Zahl am AWG, frei einstellbar zwischen $0$ und
$2\pi$. Warum also lässt sich der Crestfaktor damit nicht senken?

Er lässt sich senken. Aber es ist genau die Phase, die das
Zentrierungskriterium bereits vergeben hat. Das ist kein Zufall, und der
folgende Abschnitt rechnet es explizit aus -- weil die Größenordnung des
Tauschs die eigentliche Entscheidung ist.

\paragraph{Abzählen der wirksamen Freiheiten.}
Von den drei Phasen $(\varphi_0,\varphi_1,\varphi_2)$ sind zwei ohne Wirkung
auf $\crest$:
\begin{enumerate}
  \item Eine \emph{globale} Phase multipliziert $A(t)$ mit einem Betrag-Eins-Faktor.
        Weder $\max|A|$ noch $\mathrm{rms}|A|$ merken das.
  \item Eine \emph{Rampe} ist nach dem vorigen Abschnitt eine reine
        Zeitverschiebung der Einhüllenden. Am Arbeitspunkt mit
        $d=\SI{0.185}{\mega\hertz}$ entspricht die Steigung
        $\Delta\varphi=12{,}5^\circ$ dem Versatz
        $\tau=\Delta\varphi/(2\pi d)=\SI{0.188}{\micro\second}$: Der Satz
        $\varphi_x=[\,0;\,12{,}5;\,25\,]$ ist physikalisch \emph{derselbe}
        wie $[\,0;\,0;\,0\,]$, nur \SI{0.19}{\micro\second} anders
        getriggert. Als Stellschraube ist die Rampe damit redundant zu
        $t_0$ -- sie ist keine zusätzliche Freiheit, sondern eine zweite
        Schreibweise für eine schon vorhandene.
\end{enumerate}
Übrig bleibt \emph{ein} crestwirksamer Parameter, die Abweichung der
mittleren Phase von ihrem Rampenwert:
\begin{equation}
  \boxed{\;\psi \;=\; \varphi_1 - \frac{\varphi_0+\varphi_2}{2}\;}
  \label{eq:psidef}
\end{equation}
Man nennt $\psi$ die \emph{Krümmung} des Phasenprofils: Null heißt, die drei
Phasen liegen auf einer Geraden über dem Tonindex; $\psi\neq0$ biegt den
mittleren Ton aus dieser Geraden heraus.

\paragraph{Der Crestfaktor in geschlossener Form.}
Mit den Amplituden $a_0=a_2=\varrho$, $a_1=1$ und dem Tonabstand $d$ ist die
Einhüllende im mitgeführten Bezugssystem (Trägerfrequenz $\bar f = f_1$)
\begin{equation}
  A(t) = \varrho\,e^{\mathrm{i}(\varphi_0-2\pi d t)} + e^{\mathrm{i}\varphi_1}
       + \varrho\,e^{\mathrm{i}(\varphi_2+2\pi d t)} .
\end{equation}
Zieht man $e^{\mathrm{i}\varphi_1}$ heraus und schreibt die beiden
Außenphasen als Mittelwert plus Differenz,
\begin{equation}
  \varphi_0-\varphi_1 = -\psi+\delta, \qquad
  \varphi_2-\varphi_1 = -\psi-\delta, \qquad
  \delta = \tfrac{1}{2}(\varphi_0-\varphi_2),
\end{equation}
so fassen sich die beiden Außentöne zu einem einzigen Kosinus zusammen:
\begin{align}
  A(t)\,e^{-\mathrm{i}\varphi_1}
  &= \varrho\,e^{\mathrm{i}(-\psi+\delta-2\pi d t)} + 1 + \varrho\,e^{\mathrm{i}(-\psi-\delta+2\pi d t)}\\
  &= \varrho\,e^{-\mathrm{i}\psi}\Big[e^{\mathrm{i}(\delta-2\pi d t)}+e^{-\mathrm{i}(\delta-2\pi d t)}\Big] + 1\\
  &= 2\varrho\,e^{-\mathrm{i}\psi}\cos\Theta + 1,
     \qquad \Theta = 2\pi d\,(t-t_\delta),\quad t_\delta=\frac{\delta}{2\pi d} .
     \label{eq:n3env}
\end{align}
Hier steht die ganze Struktur der Achse auf einer Zeile. $\delta$ -- die
Rampe -- taucht ausschließlich als Zeitversatz $t_\delta$ auf und kann den
Betrag nicht ändern. $\varphi_1$ ist der herausgezogene globale Faktor.
Wirksam ist allein $\psi$, und zwar als \emph{Winkel zwischen} dem
Außentonpaar und dem mittleren Ton. Der Betrag folgt sofort:
\begin{equation}
  |A|^2 = 4\varrho^2\cos^2\Theta + 4\varrho\cos\psi\,\cos\Theta + 1 .
\end{equation}
Über eine Periode durchläuft $\cos\Theta$ das ganze Intervall $[-1,1]$, das
Maximum wird bei $\cos\Theta=\operatorname{sign}(\cos\psi)$ angenommen:
\begin{empheq}[box=\fbox]{align}
  \max_t|A| &= \sqrt{4\varrho^2 + 4\varrho\,|\cos\psi| + 1}, &
  \mathrm{rms}|A| &= \sqrt{2\varrho^2+1} \;\;\text{(fest)},\\[2pt]
  \crest_x(\psi) &= \sqrt{2}\;\frac{\sqrt{4\varrho^2+4\varrho|\cos\psi|+1}}{\sqrt{2\varrho^2+1}} . &&
  \label{eq:crestpsi}
\end{empheq}
Der Nenner ist nach Abschnitt~\ref{sec:rmsphase} phasenunabhängig -- nur der
Zähler bewegt sich. Er ist maximal bei $|\cos\psi|=1$ und minimal bei
$\cos\psi=0$:
\begin{align}
  \psi = 0 \text{ oder } \pi:&\quad \max|A| = 2\varrho+1 = \textstyle\sum_n a_n
    &&\text{(alle Zeiger ausgerichtet)},\\
  \psi = \pm\tfrac{\pi}{2}:&\quad \max|A| = \sqrt{4\varrho^2+1}
    &&\text{(mittlerer Ton in Quadratur)} .
\end{align}
Mit $\varrho=\sqrt{r_x}=\sqrt{0{,}97}=0{,}98489$ und $\mathrm{rms}|A|=1{,}71464$:

\begin{center}
\begin{tabular}{lcccccc}
  \toprule
  $\psi$ & $0^\circ$ & $45^\circ$ & $60^\circ$ & $\mathbf{90^\circ}$ & $135^\circ$ & $180^\circ$\\
  \midrule
  $\crest_x$ & $2{,}4494$ & $2{,}2836$ & $2{,}1586$ & $\mathbf{1{,}8220}$ & $2{,}2836$ & $2{,}4494$\\
  \bottomrule
\end{tabular}
\end{center}

\noindent Gl.~\eqref{eq:crestpsi} und die direkte numerische Auswertung von
$\crest$ auf dem Zeitraster stimmen auf sechs Stellen überein.

\paragraph{Und jetzt die Pointe.}
Das Zentrierungskriterium Gl.~\eqref{eq:centering} lautet bei $N=3$
ausgeschrieben $\varphi_0+\varphi_2 = 2\varphi_1$, in der Sprache von
Gl.~\eqref{eq:psidef} also schlicht
\begin{equation}
  \psi = 0 \quad(\text{oder } \pi;\ \text{beide Zweige liefern}\ |\cos\psi|=1) .
\end{equation}

\medskip
\noindent\fbox{\parbox{0.96\textwidth}{\textbf{Der einzige crestwirksame
Parameter, den drei Töne besitzen, ist genau derjenige, den das
Zentrierungskriterium festlegt -- und es stellt ihn auf den Wert, bei dem
$\max|A|$ sein Maximum $\sum_n a_n$ annimmt.} Dass $\crest_x$ bei drei Tönen
„immer maximal`` ist, ist keine Eigenheit der Amplituden, sondern die
Zentrierung selbst.}}
\medskip

\paragraph{Was das Aufgeben der Zentrierung bringt -- und kostet.}
Damit ist die Frage beantwortbar, ob sich der Tausch lohnt. Dazu wurde für
feste Krümmung $\psi$ jeweils der Rest neu optimiert: die vier
$y$-Tonphasen und der Trigger $t_0$, Zielgröße $\Uw$, unter der
Nebenbedingung $\langle\theta\rangle \ge$ Zeitmittel (nicht in eine Delle
triggern). Ergebnis in Tabelle~\ref{tab:psitrade}.

\begin{table}[htb]
  \centering
  \caption{Preis der Krümmung. $\crest_y$ ist hier nicht mitbeschränkt und
  läuft in allen drei Fällen auf $\approx 2{,}80$; die relative mittlere
  Leistung ist $\propto 1/(\crest_x\crest_y)^2$ nach
  Gl.~\eqref{eq:powerbudget}.}
  \label{tab:psitrade}
  \begin{tabular}{lccccc}
    \toprule
    $\psi$ & $\crest_x$ & $(\crest_x\crest_y)^2$
           & $P_\mathrm{mittel}$ rel. & $\Uw$ & Schwerpunktversatz\\
    \midrule
    $0^\circ$ (zentriert) & $2{,}4494$ & $47{,}9$ & $1{,}00$ & $\mathbf{0{,}94}\,\%$ & $0$\,nm\\
    $45^\circ$            & $2{,}2836$ & $40{,}7$ & $1{,}18$ & $1{,}27\,\%$ & $0$\,nm\\
    $90^\circ$            & $\mathbf{1{,}8220}$ & $25{,}9$ & $\mathbf{1{,}85}$ & $3{,}88\,\%$ & $0$\,nm\\
    \bottomrule
  \end{tabular}
\end{table}

Die Krümmung $\psi=90^\circ$ ist also real \emph{Faktor $1{,}85$ mehr
mittlere HF-Leistung} wert, $\SI{+2.7}{\decibel}$ auf der $x$-Achse. Das ist
kein Rundungseffekt, sondern derselbe Hebel, der in
Abschnitt~\ref{sec:crestback} die $y$-Achse \SI{65}{\percent} wert war.

\paragraph{Wo genau der Schaden auftritt.}
Bemerkenswert ist die letzte Spalte: Der Schwerpunktversatz bleibt in allen
drei Zeilen null. Das Kriterium ist für einen verschwindenden Schwerpunkt
also \emph{hinreichend, aber nicht notwendig} -- $t_0$ und die
$y$-Phasen können das ungerade Moment der \emph{zeitintegrierten}
Pulsfläche auch dann noch zu Null bringen, wenn die Symmetrie zu keinem
einzelnen Zeitpunkt gilt. Was verlorengeht, ist nicht die Lage, sondern die
Ruhe: Bei $\psi=0$ steht das Interferenzmuster über den ganzen Puls still,
bei $\psi\neq0$ atmet es und verschiebt sich hin und her, und nur die Bilanz
über die Pulsdauer hebt sich auf. Übrig bleibt der \emph{gerade} Anteil --
die Breite der $\theta$-Verteilung über die Wolke --, und der wächst:
$\Uw$ von $0{,}94\,\%$ auf $3{,}88\,\%$, ein Faktor vier.

Übersetzt auf einen $\pi$-Puls: Eine relative Streuung $\Uw$ der Pulsfläche
gibt eine Winkelstreuung $\pi\,\Uw$ und damit eine Infidelität
\begin{equation}
  1-F \;\approx\; \left(\frac{\pi\,\Uw}{2}\right)^{\!2}
  \;=\;
  \begin{cases}
    2{,}2\cdot10^{-4} & (\Uw = 0{,}94\,\%),\\[2pt]
    3{,}7\cdot10^{-3} & (\Uw = 3{,}88\,\%).
  \end{cases}
\end{equation}

\paragraph{Entscheidungsregel.}
Der Tausch ist damit sauber bezifferbar und hat keine allgemeingültige
Antwort, sondern hängt daran, welche der beiden Grenzen zuerst greift:

\begin{itemize}
  \item \textbf{Klemmt der Verstärker} -- reicht die erreichbare Beugung
    nach Abschnitt~\ref{sec:realsetup} nicht für die gewünschte
    Rabifrequenz --, so sind $\psi=90^\circ$ und $\SI{0.4}{\percent}$
    Infidelität aus der Ortsverteilung ein guter Handel: Faktor $1{,}85$ in
    der Leistung ist über $\eta=\sin^2(\kappa\sqrt{P})$ ein erheblicher
    Sprung, und $3{,}7\cdot10^{-3}$ liegt unter dem, was Streuung und
    Lichtverschiebung ohnehin beitragen.
  \item \textbf{Klemmt die Fidelity}, so bleibt $\psi=0$ gesetzt, und die
    Leistung wird dort geholt, wo sie noch frei ist: auf der $y$-Achse über
    den Parameter $b$ (Abschnitt~\ref{sec:crestback}), und ab $N_x\ge4$
    Tönen auch auf der $x$-Achse, wo nach Tabelle~\ref{tab:dims} das
    Kriterium fast gratis wird.
\end{itemize}

\noindent Der zweite Punkt ist die eigentliche Empfehlung: Die Enge bei
$N=3$ ist eine Eigenschaft der \emph{Tonzahl}, nicht der Physik. Ein
vierter $x$-Ton löst den Konflikt, statt ihn auszuhandeln.

\subsection{Wofür der Crestfaktor am Ende entscheidend ist}

Damit lässt sich die Rolle des Crestfaktors scharf fassen. Man vergleiche
\num{220} zufällige Phasensätze \emph{aus} der zentrierten Familie mit
\num{220} völlig freien Sätzen, jeweils bei bestmöglichem $t_0$ unter der
Bedingung $A(t_0)\ge\langle A\rangle$:

\begin{table}[htb]
  \centering
  \caption{Die Zentrierung erledigt die Homogenität; innerhalb der Familie
  unterscheidet sie die Mitglieder kaum noch.}
  \label{tab:decisive}
  \begin{tabular}{lcc}
    \toprule
     & aus der Familie & völlig frei\\
    \midrule
    $\Uw$ Median                 & \SI{0.28}{\percent} & \SI{4.15}{\percent}\\
    $\Uw$ 10--90-Perzentil       & \SIrange{0.12}{0.55}{\percent} & \SIrange{1.49}{12.45}{\percent}\\
    $\crest_y$ Spanne            & $2{,}197 \dots 2{,}826$ & --\\
    $A(t_0)/\langle A\rangle$    & $1{,}00 \dots 3{,}90$ & --\\
    \bottomrule
  \end{tabular}
\end{table}

Zwei Dinge stehen hier nebeneinander:

\begin{enumerate}\itemsep0.3em
  \item \textbf{Das Zentrierungskriterium leistet die Physik.} Es drückt $\Uw$
        im Median von \SI{4.15}{\percent} auf \SI{0.28}{\percent} -- Faktor
        \num{15} -- und garantiert dabei den Versatz exakt null. Das ist der
        Schritt, der über die Qualität der Anregung entscheidet.
  \item \textbf{Innerhalb der Familie ist $\Uw$ praktisch erledigt.} Die
        \numrange{0.12}{0.55} Prozent des 10--90-Perzentils liegen allesamt
        weit unter jeder anderen Fehlerquelle des Experiments; $\Uw$
        unterscheidet die Mitglieder nicht mehr sinnvoll. Was sie
        unterscheidet, ist der Crestfaktor -- Spanne $2{,}197$ bis $2{,}826$,
        nach Gl.~\eqref{eq:intensity} ein Faktor $1{,}66$ in der zulässigen
        mittleren Leistung -- und die erreichbare Pulsfläche.
\end{enumerate}

\medskip
\noindent\textbf{Der Crestfaktor ist also nicht das primäre Optimierungsziel,
sondern das Auswahlkriterium \emph{innerhalb} der durch die Physik bereits
festgelegten Familie.} Erst wird die Zentrierung als harte Bedingung gesetzt,
dann wird die verbleibende Freiheit -- bei $3\times4$ ausschließlich die eine
Dimension der $y$-Achse -- zur Entlastung der HF-Kette genutzt. In dieser
Reihenfolge ist er optimierungsentscheidend; als eigenständiges Ziel wäre er
es nicht, denn die crest-minimalen Sätze (Schroeder) verletzen gerade die
Bedingung, auf die es ankommt.

\section{Quantifizierung am Arbeitspunkt}

Arbeitspunkt: $3\times4$ Töne, Airy-Profil, Waist \SI{1.04}{\micro\meter},
$\Delta f = \SI{0.37}{\mega\hertz}$, $r_x/r_y = 0{,}97/1{,}16$,
$f_1/f_2 = 75/\SI{750}{\milli\meter}$, Pulslänge $T_p=\SI{1}{\micro\second}$,
Atom bei \SI{17}{\micro\kelvin} und $\nu_r=\SI{60.4}{\kilo\hertz}$, also
$\sigma=\SI{108}{\nano\meter}$.

\begin{table}[htb]
  \centering
  \caption{Vergleich dreier Phasensätze. $A(t_0)/\langle A\rangle$ ist die
  atomgewichtete Pulsfläche am Trigger, bezogen auf ihren Zeitmittelwert.}
  \label{tab:vergleich}
  \begin{tabular}{lccccc}
    \toprule
    Phasensatz & $\crest_x$ & $\crest_y$ & $\Uw$ & Versatz & $A(t_0)/\langle A\rangle$\\
    \midrule
    alle Phasen $0$ & $2{,}449$ & $2{,}826$ & \SI{0.18}{\percent} & \SI{0.0}{\nano\meter} & $1{,}35$\\
    \textbf{gewählter Punkt} & $2{,}449$ & $\mathbf{2{,}196}$ & \SI{0.30}{\percent} & \SI{0.0}{\nano\meter} & $1{,}92$\\
    frühere Wahl    & $2{,}449$ & $2{,}485$ & \SI{0.21}{\percent} & \SI{0.0}{\nano\meter} & $1{,}15$\\
    Schroeder       & $2{,}159$ & $2{,}000$ & \SI{17.9}{\percent} & \SI{18.9}{\nano\meter} & $0{,}86$\\
    \bottomrule
  \end{tabular}
\end{table}

Der gewählte Punkt ist
\begin{equation}
  \varphi_x=[\,0;\;12{,}5;\;25{,}0\,]^\circ,\qquad
  \varphi_y=[\,0;\;98{,}0;\;16{,}0;\;114{,}0\,]^\circ,\qquad
  t_0=\SI{15.306}{\micro\second}.
  \label{eq:workingpoint}
\end{equation}
Er erfüllt Gl.~\eqref{eq:centering} auf beiden Achsen -- Spiegelsummen
$25{,}0^\circ$ bzw.\ $114{,}0^\circ$, Abweichungen vom Mittenwinkel
$\chi_x=[-12{,}5;\,0;\,+12{,}5]^\circ$ und
$\chi_y=[-57;\,+41;\,-41;\,+57]^\circ$, beide antisymmetrisch -- und hält die
frequenzentarteten Spots in Quadratur ($|\cos\Delta\varphi| = 0{,}018$).
Nachgerechnet: $\Uw=\SI{0.304}{\percent}$, Versatz \SI{0.0000}{\nano\meter},
$A(t_0)=1{,}92\,\langle A\rangle$.

Die $y$-Achse schöpft dabei die einzige verbliebene Freiheit aus:
$\crest_y = 2{,}196$ ist das Minimum, das Gl.~\eqref{eq:centering} bei vier
Tönen zulässt. Eine frühere Wahl,
$\varphi_x=[\,0;\,175{,}6;\,351{,}2\,]^\circ$ und
$\varphi_y=[\,0;\,10{,}5;\,250{,}7;\,261{,}2\,]^\circ$, ist ebenso exakt
zentriert, lässt mit $\crest_y=2{,}485$ aber \SI{28}{\percent} HF-Leistung
liegen; sie steht in den Tabellen zum Vergleich.

\paragraph{Was die Phasenwahl an HF-Reserve bringt.}
Gegenüber $\varphi=0$ bleibt $\crest_x$ unverändert -- es kann nach
Abschnitt~\ref{sec:crestback} gar nicht anders sein --, während $\crest_y$ von
$2{,}826$ auf $2{,}196$ fällt. Nach Gl.~\eqref{eq:intensity} entspricht das
\begin{equation}
  \frac{I}{I_{\varphi=0}}
  = \left(\frac{2{,}826}{2{,}196}\right)^{2} = 1{,}66
  \qquad\widehat{=}\qquad \SI{2.2}{\deci\bel}
\end{equation}
mehr zulässige mittlere Leistung. Schroeder-Phasen kämen auf
\SI{4.10}{\deci\bel}, verletzen aber Gl.~\eqref{eq:centering}: Der Fleck sitzt
dann $0{,}18\,\sigma$ neben dem Atom und $\Uw$ steigt um zwei
Größenordnungen.

\paragraph{Der Preis der Zentrierung.}
Die Differenz ist die eigentlich zu begründende Zahl:
\begin{equation}
  \SI{4.10}{\deci\bel} - \SI{2.2}{\deci\bel} \;\approx\; \SI{1.9}{\deci\bel}
\end{equation}
Verstärkerreserve kostet die exakte Zentrierung gegenüber der reinen
Crest-Minimierung. Die Begründung lautet damit: \emph{Die
Zentrierungsbedingung kostet rund \SI{2}{\deci\bel} gegenüber dem
crest-optimalen Satz, und das ist gerechtfertigt, weil andernfalls der
Interferenzfleck um $0{,}18\,\sigma$ neben dem Atom steht und $\Uw$ von
\SI{0.3}{\percent} auf \SI{18}{\percent} springt.}

\section{Abgrenzung: was der Crestfaktor nicht ist}\label{sec:abgrenzung}

Der Crestfaktor wird leicht mit dem Verhältnis von maximaler zu gemittelter
\emph{Intensität} verwechselt, zumal beide am Arbeitspunkt zufällig
ähnliche Zahlenwerte annehmen. Sie sind verschieden und nicht einmal monoton
gekoppelt, wie Tabelle~\ref{tab:abgrenzung} zeigt.

\begin{table}[htb]
  \centering
  \caption{HF-Crestfaktor gegen optische Spitzenwerte. $I_W$ ist die
  atomgewichtete Intensität, „Profil`` das Maximum über Ort und Zeit,
  bezogen auf das Maximum des Zeitmittels.}
  \label{tab:abgrenzung}
  \begin{tabular}{lccc}
    \toprule
    Phasensatz & $\crest_x/\crest_y$ & $\max_t I_W/\langle I_W\rangle$ & Profil\\
    \midrule
    alle Phasen $0$ & $2{,}449/2{,}826$ & $4{,}44$ & $4{,}18$\\
    gewählter Punkt & $2{,}449/2{,}485$ & $2{,}92$ & $4{,}04$\\
    Schroeder       & $2{,}159/2{,}000$ & $2{,}86$ & $\mathbf{5{,}17}$\\
    \bottomrule
  \end{tabular}
\end{table}

Schroeder hat den niedrigsten HF-Crestfaktor und zugleich den höchsten
optischen Spitzenwert im Profil -- eine Identität der beiden Größen ist damit
ausgeschlossen. Im Code wird der Crestfaktor ausschließlich aus $f_n$,
$\varphi_n$ und $a_n$ berechnet; Feld, Ortsgitter und Intensität gehen nicht
ein.

Ein Zusammenhang besteht lediglich im Grenzfall vollständig überlappender
Spots. Dort wäre $I \propto |A_x|^2|A_y|^2$ und folglich
\begin{equation}
  \frac{\max_t I}{\langle I\rangle_t}
  \;\approx\; \frac{\crest_x^{2}}{2}\cdot\frac{\crest_y^{2}}{2},
\end{equation}
für den gewählten Punkt also $9{,}26$; die direkte Auswertung dieses
Grenzfalls ergibt $6{,}71$ (die Abweichung rührt daher, dass die Töne
kommensurabel sind und die beiden Einhüllenden ihr Maximum nicht unabhängig
voneinander erreichen). In der realen Geometrie sind die Spots getrennt, am
Atomort trägt nur die Teilüberlappung bei, und der Wert sinkt auf $2{,}92$.

Beide Größen adressieren zudem verschiedene Budgets: Der HF-Crestfaktor
begrenzt Verstärkerkompression und Spitzenlast am AOD-Wandler, die optische
Spitzenintensität dagegen Lichtverschiebung und Streuung am Atom.

\section{Fazit}

Der Crestfaktor ist keine physikalische Einschränkung der Tonphasen -- jeder
Satz ist einstellbar und wird von der Simulation korrekt behandelt. Er ist
eine Aussage über die Hardware: Gl.~\eqref{eq:powerbudget} und
\eqref{eq:intensity} übersetzen ihn in Licht am Atom. Als Nebenbedingung ist
er dadurch gerechtfertigt, dass er die Voraussetzung der Rechnung -- lineare
Beugung ohne Intermodulation -- überhaupt erst sichert.

Die Argumentationskette für die Parameterwahl lautet damit:
\begin{enumerate}\itemsep0.2em
  \item Nur Phasendifferenzen gehen ein; eine globale Phase je Achse und eine
        gemeinsame, in der Zeit abgestimmte Rampe auf beiden Achsen sind
        Eichfreiheiten -- letztere ist exakt eine Verschiebung von $t_0$
        (numerisch auf $3\cdot10^{-11}$ geprüft).
  \item Physikalisches Kriterium: gleichmäßige Beleuchtung des Atoms, also
        minimales $\Uw$. Gl.~\eqref{eq:centering} erfüllt die Zentrierung
        exakt und per Konstruktion statt per Anpassung.
  \item Hardwarekriterium: $\crest_x$ und $\crest_y$ unterhalb der
        gemessenen Kompressionsgrenze der Kette.
  \item Die verbleibende Freiheit -- bei $3\times4$ ausschließlich auf der
        $y$-Achse -- wird zur Absenkung des Crestfaktors genutzt.
\end{enumerate}

\section{Die Rechnung am realen Aufbau}\label{sec:realsetup}

Dieser Abschnitt rechnet alles Vorangegangene an der tatsächlichen Hardware
durch, Schritt für Schritt.

\subsection{Drei Grenzen, drei Geräte}

Man verwechselt leicht, was wovon begrenzt wird. Es sind drei unabhängige
Budgets:

\begin{center}
\begin{tabular}{lll}
  \toprule
  Grenze & gehört zu & hängt vom Crestfaktor ab?\\
  \midrule
  \SI{2}{\watt} mittlere HF-Leistung & AOD-Wandler (Wärme) & \textbf{nein}\\
  Kompression & Verstärker (Spitzenspannung) & \textbf{ja}, $\propto 1/\crest^2$\\
  Beugungssättigung $\sin^2$ & AOD (Leistung \emph{pro Ton}) & nein\\
  \bottomrule
\end{tabular}
\end{center}

Die \SI{2}{\watt} des AOD sind ein \emph{Mittelwert} über die Zeit -- der
Wandler wird warm, und wie die Leistung innerhalb einer Mikrosekunde verteilt
ist, ist ihm gleich. Der Verstärker dagegen sieht die Momentanspannung; ihn
interessiert nur die Spitze.

\subsection{Die Daten}

\begin{center}
\begin{tabular}{ll}
  \toprule
  \multicolumn{2}{l}{\textbf{Verstärker Mini-Circuits ZHL-03-5WF+}}\\
  Frequenzbereich & \SIrange{60}{300}{\mega\hertz} (\SI{100}{\mega\hertz} liegt darin)\\
  Verstärkung & \SI{30}{\deci\bel}\\
  $P_{1\mathrm{dB}}$ & $+\SI{36}{dBm}$\\
  $P_\mathrm{sat}$ & $+\SI{39}{dBm}$\\
  OIP3 & $+\SI{49}{dBm}$\\
  \midrule
  \multicolumn{2}{l}{\textbf{AOD AA DTSXY-400-800, S/N 3001-O211386}}\\
  max.\ HF-Leistung je Achse & \SI{2}{\watt} (Wandlerschaden)\\
  Messpunkt im Testsheet & \SI{1.3}{\watt} je Achse\\
  Effizienz dort, Ordnung „1-1`` & $>\SI{58}{\percent}$ (Spez.\ $>\SI{40}{\percent}$)\\
  Mittenfrequenz bei \SI{795}{\nano\meter} & \SI{103}{\mega\hertz} $\pm$ \SI{4}{\mega\hertz}\\
  Bandbreite & \SI{36}{\mega\hertz}\\
  \bottomrule
\end{tabular}
\end{center}

\paragraph{dBm in Watt.}
Die Verstärkerangaben stehen in dBm, also Dezibel bezogen auf
\SI{1}{\milli\watt}:
\begin{equation}
  P[\mathrm{W}] \;=\; \frac{1}{1000}\cdot 10^{\,P[\mathrm{dBm}]/10} .
\end{equation}
Damit
\begin{align}
  P_{1\mathrm{dB}} &= \tfrac{1}{1000}\cdot10^{36/10} = \tfrac{1}{1000}\cdot 3981 = \SI{3.98}{\watt},\\
  P_\mathrm{sat}   &= \tfrac{1}{1000}\cdot10^{39/10} = \tfrac{1}{1000}\cdot 7943 = \SI{7.94}{\watt}.
\end{align}
Der \emph{Kompressionspunkt} $P_{1\mathrm{dB}}$ ist die Ausgangsleistung, bei
der die Verstärkung um \SI{1}{\deci\bel} eingebrochen ist -- ab dort ist die
Kette nicht mehr linear. Genau die Linearität setzt die Feldrechnung voraus,
also ist $P_{1\mathrm{dB}}$ die richtige Grenze, nicht $P_\mathrm{sat}$.

\subsection{Schritt 1: Welche mittlere Leistung erlaubt der Verstärker?}

Entscheidend ist: Die Grenze gilt für die \emph{Spitze}, nicht für den
Mittelwert. Nach Gl.~\eqref{eq:powerbudget} ist
$P_\mathrm{peak}=\crest^2 P_\mathrm{mittel}$, die Forderung
$P_\mathrm{peak}\le P_{1\mathrm{dB}}$ liefert also
\begin{equation}
  P_\mathrm{mittel} \;\le\; \frac{P_{1\mathrm{dB}}}{\crest^{2}} .
  \label{eq:pavglimit}
\end{equation}
Eingesetzt, mit den Crestfaktoren des gewählten Arbeitspunkts:
\begin{align}
  \text{$x$-Achse:}\quad &\crest_x = 2{,}449
  &&\crest_x^2 = 6{,}00
  &&P_\mathrm{mittel,x} \le \frac{3{,}98}{6{,}00} = \SI{0.66}{\watt},\\
  \text{$y$-Achse:}\quad &\crest_y = 2{,}196
  &&\crest_y^2 = 4{,}82
  &&P_\mathrm{mittel,y} \le \frac{3{,}98}{4{,}82} = \SI{0.83}{\watt}.
\end{align}

\medskip
\noindent\fbox{\parbox{0.96\textwidth}{\textbf{Das ist das Ergebnis, auf das
es ankommt:} Erlaubt sind \SI{0.66}{\watt} und \SI{0.83}{\watt} -- höchstens
\SI{0.4}{} der AOD-Grenze von \SI{2}{\watt} und deutlich unter dem Messpunkt
von \SI{1.3}{\watt}. \emph{Der AOD ist nicht das Limit, der Verstärker ist
es.} Der Wandler wird bei diesem Betrieb nie warm.}}

\subsection{Schritt 2: Gegenprobe}

Umgekehrt gefragt: Welche Spitzenleistung bräuchte man, um die AOD-Grenzen
auszufahren? Maßgeblich ist die $x$-Achse mit $\crest_x^2=6{,}00$:
\begin{align}
  P_\mathrm{mittel} = \SI{1.3}{\watt} &\;\rightarrow\; P_\mathrm{peak} = 6{,}00\cdot1{,}3 = \SI{7.8}{\watt} = +\SI{38.9}{dBm},\\
  P_\mathrm{mittel} = \SI{2.0}{\watt} &\;\rightarrow\; P_\mathrm{peak} = 6{,}00\cdot2{,}0 = \SI{12.0}{\watt} = +\SI{40.8}{dBm}.
\end{align}
Die \SI{1.3}{\watt} landen exakt auf $P_\mathrm{sat}$ -- der Verstärker
begrenzt dort bereits hart, von Linearität keine Rede. Die vollen
\SI{2}{\watt} lägen \SI{1.9}{\deci\bel} \emph{über} $P_\mathrm{sat}$ und sind
mit dieser Kette gar nicht erreichbar.

\subsection{Schritt 3: Von der HF-Leistung zur Beugung}\label{sec:beugung}

Jetzt ist die Frage zu beantworten, wie viel Licht die erlaubte HF-Leistung
tatsächlich in das Profil bringt. Drei Dinge sind dafür zu klären, und nur
das dritte ist wirklich heikel:
\begin{enumerate}\itemsep0.15em
  \item Wie macht \emph{ein} Gitter aus HF-Leistung Beugung? (Abschnitt unten,
        Ergebnis $\eta=\sin^2(\kappa\sqrt{P})$.)
  \item Woher kommt die Zahl $\kappa$? (Aus dem Testsheet, nicht aus einer Rechnung.)
  \item Was ändert sich, wenn $N$ Gitter \emph{gleichzeitig} stehen?
        (Die Summe wandert unter die Wurzel.)
\end{enumerate}

\paragraph{Woher das \boldmath$\sin^2$ kommt, und warum \boldmath$\sqrt{P}$.}
Die HF-Leistung erzeugt am Piezowandler eine akustische Welle. Deren
Dehnungsamplitude $S$ hängt mit der akustischen Leistung über
$P_\mathrm{ak}\propto S^2$ zusammen -- Leistung ist quadratisch in der
Amplitude, wie überall. Über den photoelastischen Effekt wird daraus eine
Brechungsindexmodulation
\begin{equation}
  \Delta n = -\tfrac{1}{2}\,n^3 p\,S \;\;\propto\;\; \sqrt{P_\mathrm{ak}} ,
  \label{eq:photoelast}
\end{equation}
also ein Phasengitter, dessen \emph{Amplitude} mit der Wurzel der Leistung
wächst. \textbf{Das ist der Ursprung der Wurzel} -- sie steckt nicht in der
Beugung, sondern in der Umrechnung von Leistung auf Amplitude.

Im Bragg-Regime tauschen einfallende und gebeugte Welle entlang der
Wechselwirkungsstrecke $z$ periodisch Amplitude aus. Es sind zwei gekoppelte
Wellen mit Kopplungskonstante $\propto\Delta n$; mit $A_\mathrm{u}(0)=1$, $A_\mathrm{d}(0)=0$
lautet die Lösung $A_\mathrm{u}=\cos(\nu z/L)$, $A_\mathrm{d}=-\mathrm{i}\sin(\nu z/L)$ und
damit am Ausgang
\begin{equation}
  \eta = |A_\mathrm{d}|^2 = \sin^2\nu , \qquad
  \nu = \frac{\pi\,\Delta n\,L}{\lambda\cos\theta_\mathrm{B}}
      = \frac{\pi}{\lambda}\sqrt{\frac{M_2\,L\,P_\mathrm{ak}}{2H}} ,
\end{equation}
mit der Wechselwirkungslänge $L$, der Wandlerhöhe $H$ und der Gütezahl $M_2$
des Kristalls. Zusammengefasst zu einer einzigen Konstanten:
\begin{equation}
  \boxed{\;\eta(P) = \sin^2\!\bigl(\kappa\sqrt{P}\bigr)\;}
  \qquad
  \kappa = \frac{\pi}{\lambda}\sqrt{\frac{M_2 L}{2H}}\cdot\sqrt{\zeta} ,
  \label{eq:etasingle}
\end{equation}
wobei $\zeta$ den Wirkungsgrad der elektroakustischen Wandlung und der
HF-Anpassung enthält, denn $P$ soll die \emph{elektrische} Leistung sein.

\medskip
\noindent In $\kappa$ stecken damit vier Größen -- $M_2$, das
Aspektverhältnis $L/H$, die Wandlung $\zeta$ und die Anpassung --, von denen
keine einzeln mit brauchbarer Genauigkeit bekannt ist. Genau deshalb wird
$\kappa$ \emph{gemessen} und nicht gerechnet: Eine einzige Effizienzmessung
bei einer bekannten Leistung ersetzt alle vier.

\paragraph{Der Ankerpunkt aus dem Testsheet.}
Das Testsheet nennt \SI{58}{\percent} in der Ordnung „1-1`` bei
\SI{1.3}{\watt} je Achse, Einzelton. „1-1`` heißt einmal in $x$ \emph{und}
einmal in $y$ gebeugt; die \SI{58}{\percent} sind also bereits das
\emph{Produkt} beider Achsen. Unter der Annahme, dass die beiden Achsen
desselben Bauteils gleich gut sind -- sie teilen Kristall, Wandlerprozess und
Testsheet --, folgt je Achse
\begin{equation}
  \eta_\mathrm{Achse} = \sqrt{0{,}58} = 0{,}762 \quad\text{bei } \SI{1.3}{\watt},
\end{equation}
und Gl.~\eqref{eq:etasingle} nach $\kappa$ aufgelöst:
\begin{equation}
  \kappa = \frac{\arcsin\sqrt{0{,}762}}{\sqrt{1{,}3}}
         = \frac{1{,}061}{1{,}140} = \SI{0.930}{\per\sqrt\watt}.
  \label{eq:kappa}
\end{equation}
Zwei Vorbehalte dazu: Das Testsheet schreibt „$>\SI{58}{\percent}$``, der
Wert ist also eine \emph{untere} Schranke, und er stammt aus einer
Einzeltonmessung. Beides macht die folgenden Zahlen eher konservativ.

\paragraph{Woher die gekoppelten Gleichungen kommen.}
Gl.~\eqref{eq:etasingle} wurde oben als Ergebnis zitiert; für den
Mehrtonfall genügt das nicht, weil dort gerade die \emph{Struktur} der
Gleichungen die Antwort trägt. Also von vorn.

\medskip
\noindent\textbf{Ausgangspunkt.} Im Kristall gilt die skalare
Helmholtz-Gleichung
\begin{equation}
  \nabla^2 E + k_\mathrm{vac}^2\,n^2(\mathbf{r},t)\,E = 0 ,
  \qquad k_\mathrm{vac} = \frac{\omega_0}{c} ,
\end{equation}
und die Schallwellen modulieren den Index. Jeder Ton $n$ ist eine laufende
Welle mit Wellenvektor $\mathbf{K}_n$ ($|\mathbf{K}_n| = \Omega_n/v$,
$v$ = Schallgeschwindigkeit), Kreisfrequenz $\Omega_n = 2\pi f_n$ und
\emph{seiner} Phase $\varphi_n$:
\begin{equation}
  n(\mathbf{r},t) = n_0 + \sum_n \Delta n_n
    \cos\!\bigl(\mathbf{K}_n\!\cdot\!\mathbf{r} - \Omega_n t + \varphi_n\bigr),
  \qquad \Delta n_n \propto \sqrt{P_n}
\end{equation}
nach Gl.~\eqref{eq:photoelast}. Für $\Delta n_n\ll n_0$ ist
$n^2 \approx n_0^2 + 2n_0\sum_n\Delta n_n\cos(\dots)$ -- der quadratische
Term ist die akustische Nichtlinearität aus Punkt~(iii) unten und wird hier
weggelassen.

\medskip
\noindent\textbf{Ansatz.} Das Feld ist die Überlagerung \emph{einer}
ungebeugten Welle und $N$ gebeugter, jede mit langsam veränderlicher
Amplitude:
\begin{equation}
  E = \underbrace{A_\mathrm{u}(z)\,
        e^{\mathrm{i}(\mathbf{k}_\mathrm{u}\cdot\mathbf{r}-\omega_0 t)}}_{
        \text{ungebeugt}}
    \;+\; \underbrace{\sum_{n=0}^{N-1} A_n(z)\,
        e^{\mathrm{i}(\mathbf{k}_n\cdot\mathbf{r}-\omega_n t)}}_{
        \text{je ein Strahl pro Ton}} ,
  \qquad
  \omega_n = \omega_0 + \Omega_n .
  \label{eq:ansatz}
\end{equation}
Alle Wellenvektoren liegen auf der Dispersionskugel
($|\mathbf{k}_n| = n_0\omega_n/c$); $\mathbf{k}_n$ zeigt in die Richtung, in
die die $n$-te Ordnung tatsächlich läuft. „Langsam veränderlich`` heißt
$|\mathrm{d}^2A/\mathrm{d}z^2| \ll |k\,\mathrm{d}A/\mathrm{d}z|$, die
zweite Ableitung fällt also weg.

\medskip
\noindent\fbox{\parbox{0.96\textwidth}{\textbf{Zur Bezeichnung.} Der
\emph{ungebeugte} Strahl heißt hier $A_\mathrm{u}$ und trägt \emph{keinen}
Zahlenindex -- er gehört zu keinem Ton. Der Tonindex $n$ läuft wie im ganzen
Text von $0$ bis $N-1$, und $A_n$ ist der von Ton $n$ erzeugte gebeugte
Strahl. Insbesondere ist $A_0$ die erste Beugungsordnung des
\emph{ersten Tons}, nicht der durchgehende Strahl. Jede Summe $\sum_n$
in diesem Abschnitt läuft über die $N$ Töne und enthält
$A_\mathrm{u}$ nie.}}

\medskip
\noindent\textbf{Einsetzen und sortieren.} Der Kosinus im Indexterm zerfällt
in $e^{+\mathrm{i}(\dots)}$ und $e^{-\mathrm{i}(\dots)}$. Multipliziert man
$A_\mathrm{u}\,e^{\mathrm{i}(\mathbf{k}_\mathrm{u}\mathbf{r}-\omega_0 t)}$ mit dem
$+$-Anteil des Tons $n$, so entsteht ein Term mit Raum-Zeit-Abhängigkeit
$e^{\mathrm{i}[(\mathbf{k}_\mathrm{u}+\mathbf{K}_n)\mathbf{r}-(\omega_0+\Omega_n)t]}$
-- in der Zeit genau $\omega_n$, im Raum aber
$\mathbf{k}_\mathrm{u}+\mathbf{K}_n$, und das ist \emph{nicht} notwendig
$\mathbf{k}_n$. Sammelt man alle Beiträge, die mit $\omega_n$ schwingen, und
bezieht sie auf den Träger $e^{\mathrm{i}\mathbf{k}_n\mathbf{r}}$ aus
Gl.~\eqref{eq:ansatz}, bleibt genau dieser Rest stehen:
\begin{empheq}[box=\fbox]{align}
  \frac{\mathrm{d}A_n}{\mathrm{d}z}
    &= -\mathrm{i}\,\kappa_n\,e^{\mathrm{i}\varphi_n}\,A_\mathrm{u}\;
       e^{\,\mathrm{i}\Delta k_n z} , \label{eq:coupledgena}\\
  \frac{\mathrm{d}A_\mathrm{u}}{\mathrm{d}z}
    &= -\mathrm{i}\sum_n \kappa_n\,e^{-\mathrm{i}\varphi_n}\,A_n\;
       e^{-\mathrm{i}\Delta k_n z} , \label{eq:coupledgenb}
\end{empheq}
mit der Kopplungskonstanten und dem Impulsfehler
\begin{equation}
  \kappa_n = \frac{\pi\,\Delta n_n}{\lambda\cos\theta_n} \;\propto\; \sqrt{P_n} ,
  \qquad
  \Delta k_n = \bigl(\mathbf{k}_\mathrm{u}+\mathbf{K}_n-\mathbf{k}_n\bigr)\cdot\hat{\mathbf{z}} .
  \label{eq:dkdef}
\end{equation}
$\Delta k_n$ misst also, um wie viel die Impulsbilanz
$\mathbf{k}_n = \mathbf{k}_\mathrm{u}+\mathbf{K}_n$ verfehlt wird, wenn
$\mathbf{k}_n$ zugleich auf der Dispersionskugel liegen soll. Genau null ist
sie nur für den einen Ton, für den der AOD exakt im Bragg-Winkel steht.

\medskip
\noindent Aus Gl.~\eqref{eq:coupledgena}--\eqref{eq:coupledgenb} liest man
drei Dinge ab, und alle drei werden gebraucht.

\medskip
\noindent\textbf{(a) Die Tonphasen lassen sich wegtransformieren.} Mit
$\tilde A_n = A_n e^{-\mathrm{i}\varphi_n}$ verschwindet $\varphi_n$ aus
beiden Gleichungen vollständig. Da $|A_n| = |\tilde A_n|$, ist die
Beugungseffizienz $\eta_n$ \emph{exakt} phasenunabhängig, und $\varphi_n$
taucht allein als Phase der gebeugten Welle wieder auf -- dort, wo sie
hingehört, nämlich in $\varphi_{nm}$ des Interferenzbildes. Das ist die
strenge Fassung des Arguments aus Abschnitt~\ref{sec:abgrenzung}: Der
Crestfaktor kann in der Beugungsrechnung gar nicht vorkommen.

\medskip
\noindent\textbf{(b) Die Phasenanpassung ist bereits hier eingebaut.} Setzt
man $\Delta k_n = 0$ für alle $n$ und transformiert die $\varphi_n$ nach (a)
weg, so wird aus Gl.~\eqref{eq:coupledgena}--\eqref{eq:coupledgenb}
unmittelbar
\begin{equation}
  \frac{\mathrm{d}A_\mathrm{u}}{\mathrm{d}z} = -\mathrm{i}\sum_n \kappa_n A_n ,
  \qquad
  \frac{\mathrm{d}A_n}{\mathrm{d}z} = -\mathrm{i}\,\kappa_n A_\mathrm{u} ,
  \qquad \kappa_n = \kappa\sqrt{P_n} .
  \label{eq:coupled}
\end{equation}
Hier und im Folgenden ist $z$ in Einheiten der Wechselwirkungslänge $L$
gemessen; der Kristallausgang liegt also bei $z=1$. Damit ist $\kappa_n$ die
über den \emph{ganzen} Kristall aufsummierte Kopplung -- genau die Größe
$\nu$ aus Gl.~\eqref{eq:etasingle}, weshalb $\kappa_n = \kappa\sqrt{P_n}$
dimensionslos ist.

\emph{Das} ist die Stelle, an der die gemeinsame Phasenanpassung
vorausgesetzt wird -- nicht später und nicht nebenbei. Alles, was folgt,
steht und fällt mit $\Delta k_n \approx 0$; wie gut das hier erfüllt ist,
beziffert Punkt~(i) unten.

\medskip
\noindent\textbf{(c) Ungleiche \boldmath$\Delta k_n$ zerstören die
Modenreduktion.} Das ist der Grund, warum es nicht genügt, die Verstimmung
als kleine Korrektur nachzureichen. Versucht man die helle Mode
mitzuführen, $B = \kappa_\mathrm{eff}^{-1}\sum_n \kappa_n A_n
e^{-\mathrm{i}\Delta k_n z}$, so liefert die Produktregel
\begin{equation}
  \frac{\mathrm{d}B}{\mathrm{d}z}
  = -\mathrm{i}\,\kappa_\mathrm{eff}A_\mathrm{u}
    \;-\;\frac{\mathrm{i}}{\kappa_\mathrm{eff}}
        \sum_n \kappa_n\,\Delta k_n\,A_n\,e^{-\mathrm{i}\Delta k_n z} .
\end{equation}
Der zweite Term ist eine \emph{neue} Kombination, mit den Gewichten
$\kappa_n\Delta k_n$ statt $\kappa_n$. Sie ist zu $B$ nur dann
proportional, wenn alle $\Delta k_n$ denselben Wert haben. Andernfalls
koppelt $B$ an sie, sie an die nächste (Gewichte $\kappa_n\Delta k_n^2$),
und so fort -- eine Hierarchie, die nicht abbricht. Zwei Sonderfälle
schließen sie:
\begin{itemize}\itemsep0.15em
  \item $\Delta k_n = 0$ für alle $n$: Gl.~\eqref{eq:coupled}, exakt lösbar.
  \item $\Delta k_n = \Delta k$ für alle $n$ (gemeinsame Verstimmung):
        Das System schließt wieder, und man erhält die verstimmte
        Zwei-Wellen-Lösung Gl.~\eqref{eq:detuned} mit
        $\nu = \kappa_\mathrm{eff}$.
\end{itemize}
Am Arbeitspunkt sind die $\Delta k_n$ zwar nicht gleich, aber alle so nahe
bei null, dass beide Bilder auf sechs Stellen dasselbe liefern -- siehe die
Tabelle in Punkt~(i).

\medskip
\noindent\textbf{Strukturprüfung.} Gl.~\eqref{eq:coupled} erhält die
Energie, wie sie muss:
\begin{equation}
  \frac{\mathrm{d}}{\mathrm{d}z}\Bigl(|A_\mathrm{u}|^2+\sum_n|A_n|^2\Bigr)
  = 2\sum_n \kappa_n\,\mathrm{Im}\bigl(A_\mathrm{u}^{*}A_n\bigr)
    - 2\sum_n \kappa_n\,\mathrm{Im}\bigl(A_\mathrm{u}^{*}A_n\bigr) = 0 .
\end{equation}
Das $-\mathrm{i}$ in \emph{beiden} Gleichungen ist dafür notwendig; ein
Vorzeichenfehler an einer Stelle würde hier sofort auffallen.

\medskip
\noindent\textbf{Wo die Frequenzverschiebung bleibt.} Die
$\Omega_n$ aus Gl.~\eqref{eq:ansatz} kürzen sich aus den
Amplitudengleichungen heraus -- für $\eta_n = |A_n|^2$ spielen sie keine
Rolle. Sie verschwinden aber nicht: Jede gebeugte Welle läuft mit
$\omega_0+\Omega_n$, und die Differenzen dieser Frequenzen sind genau das
Beating in der Fokalebene, mit Grundfrequenz $f_0$ und den Ordnungen $k$ aus
Abschnitt~\ref{sec:zentrierung}. Die Beugungsrechnung liefert die
\emph{Amplituden} der Spots, die Frequenzverschiebungen ihre
\emph{Zeitabhängigkeit}; beides zusammen ist das Feld, mit dem der Rest
dieses Textes arbeitet.

\paragraph{Mehrere Töne gleichzeitig: warum die Effizienzen nicht addiert werden.}
Hier liegt der eigentliche Stolperstein. Nahe liegt, jedem Ton seine eigene
Beugung zuzuschreiben und aufzusummieren. Das ist falsch, und der Grund ist
einfach: Die $N$ Gitter stehen gleichzeitig im selben Kristall und speisen
sich alle aus \emph{derselben} ungebeugten Welle. Sie konkurrieren um ein
gemeinsames Reservoir, statt jeweils ein eigenes anzuzapfen.

Formal ist das Gl.~\eqref{eq:coupled} -- $N+1$ gekoppelte Gleichungen, von
denen die erste zeigt, dass \emph{alle} $A_n$ an demselben $A_\mathrm{u}$ ziehen. Sie lassen sich in vier Schritten auf
zwei reduzieren, und zwar \emph{exakt}.

\medskip
\noindent\textbf{Schritt 1: die helle Mode.} Man definiert
\begin{equation}
  \kappa_\mathrm{eff} = \sqrt{\textstyle\sum_n \kappa_n^{2}} ,
  \qquad
  B = \frac{1}{\kappa_\mathrm{eff}}\sum_n \kappa_n A_n .
  \label{eq:brightmode}
\end{equation}
Dann ist einerseits $\mathrm{d}A_\mathrm{u}/\mathrm{d}z = -\mathrm{i}\,\kappa_\mathrm{eff}B$
-- das ist nur Gl.~\eqref{eq:coupled} rückwärts gelesen -- und andererseits
\begin{equation}
  \frac{\mathrm{d}B}{\mathrm{d}z}
  = \frac{1}{\kappa_\mathrm{eff}}\sum_n \kappa_n\,\frac{\mathrm{d}A_n}{\mathrm{d}z}
  = \frac{-\mathrm{i}}{\kappa_\mathrm{eff}}\sum_n \kappa_n^{2}\,A_\mathrm{u}
  = -\mathrm{i}\,\kappa_\mathrm{eff}A_\mathrm{u} .
\end{equation}
$A_\mathrm{u}$ und $B$ bilden also ein geschlossenes Zwei-Wellen-System mit
Kopplung $\kappa_\mathrm{eff}$ -- genau die Form, die schon gelöst ist.

\medskip
\noindent\textbf{Schritt 2: die dunklen Moden bleiben dunkel.} Nimmt man
irgendeine andere Linearkombination $C=\sum_n c_n A_n$ mit
$\sum_n c_n\kappa_n = 0$, so ist
\begin{equation}
  \frac{\mathrm{d}C}{\mathrm{d}z} = -\mathrm{i}\Bigl(\sum_n c_n\kappa_n\Bigr) A_\mathrm{u} = 0 ,
\end{equation}
und weil am Eingang noch nichts gebeugt ist, $C(0)=0$, bleibt $C\equiv 0$
über die ganze Strecke. Von den $N$ gebeugten Freiheitsgraden werden $N-1$
also \emph{nie} angeregt. Das ist der Grund, warum die Reduktion exakt ist
und keine Näherung.

\medskip
\noindent\textbf{Schritt 3: lösen.} Mit $A_\mathrm{u}(0)=1$, $B(0)=0$:
\begin{equation}
  A_\mathrm{u}(z) = \cos(\kappa_\mathrm{eff}z) , \qquad B(z) = -\mathrm{i}\sin(\kappa_\mathrm{eff}z) ,
  \qquad 0\le z\le 1 .
\end{equation}

\medskip
\noindent\textbf{Schritt 4: zurück auf die einzelnen Ordnungen.} Bisher ist
nur die \emph{Summe} $B$ bekannt; gesucht sind die einzelnen $A_n$. Der
Schlüssel steht schon in Gl.~\eqref{eq:coupled}: Dort treibt für \emph{alle}
$n$ dasselbe $A_\mathrm{u}$. Integriert man die zweite Gleichung mit $A_n(0)=0$, so
steht rechts für jedes $n$ dasselbe Integral:
\begin{equation}
  A_n(z) = -\mathrm{i}\,\kappa_n \underbrace{\int_0^z A_\mathrm{u}(z')\,\mathrm{d}z'}_{
           \text{für alle } n \text{ identisch}}
         = -\mathrm{i}\,\kappa_n \int_0^z \cos(\kappa_\mathrm{eff}z')\,\mathrm{d}z'
         = -\mathrm{i}\,\frac{\kappa_n}{\kappa_\mathrm{eff}}\,\sin(\kappa_\mathrm{eff}z) .
  \label{eq:anexplicit}
\end{equation}
Die Ordnungen unterscheiden sich also \emph{nur} durch den Vorfaktor
$\kappa_n$ -- ihre $z$-Abhängigkeit ist dieselbe, und ein Vergleich mit
Schritt~3 zeigt $A_n = (\kappa_n/\kappa_\mathrm{eff})\,B$, wie behauptet.

\medskip
\noindent Daraus folgt alles Weitere in drei Zeilen. \emph{Erstens} die
Gesamteffizienz am Kristallausgang $z=1$. Sie ist die Summe über die $N$
\emph{gebeugten} Strahlen -- der ungebeugte zählt nicht mit, er ist gerade
das, was übrig bleibt:
\begin{equation}
  \eta_\mathrm{ges} = \sum_{n=0}^{N-1} |A_n|^2
    = \underbrace{\frac{\sum_n\kappa_n^2}{\kappa_\mathrm{eff}^2}}_{=\,1
      \text{ nach Gl.~\eqref{eq:brightmode}}}\,\sin^2\kappa_\mathrm{eff}
    = \sin^2 \kappa_\mathrm{eff}
    \;=\; 1 - |A_\mathrm{u}|^2 .
  \label{eq:etages}
\end{equation}
Die letzte Gleichheit ist die Energiebilanz von unten gelesen: Was nicht
gebeugt wird, läuft geradeaus weiter. $\eta_\mathrm{ges}=\SI{47}{\percent}$
heißt also, dass \SI{53}{\percent} des Lichts den Kristall ungebeugt
verlassen.
\emph{Zweitens} -- und das ist die Stelle, an der die Summe unter die Wurzel
wandert -- wird $\kappa_\mathrm{eff}$ mit $\kappa_n = \kappa\sqrt{P_n}$ aus
Gl.~\eqref{eq:coupled} ausgeschrieben:
\begin{equation}
  \kappa_\mathrm{eff}
  = \sqrt{\sum_n \kappa_n^2}
  = \sqrt{\sum_n \kappa^2 P_n}
  = \kappa\,\sqrt{\sum_n P_n}
  = \kappa\sqrt{P_\mathrm{ges}} .
  \label{eq:keffP}
\end{equation}
Die Quadrate der Kopplungen sind die Leistungen -- deshalb ist die
quadratische Summe der $\kappa_n$ nichts anderes als die \emph{gewöhnliche}
Summe der Leistungen unter einer Wurzel. \emph{Drittens} die Aufteilung: Aus
Gl.~\eqref{eq:anexplicit} ist
$|A_n|^2 = (\kappa_n^2/\kappa_\mathrm{eff}^2)\sin^2\kappa_\mathrm{eff}$, und
$\kappa_n^2/\kappa_\mathrm{eff}^2 = P_n/P_\mathrm{ges}$. Zusammen:
\begin{empheq}[box=\fbox]{equation}
  \eta_\mathrm{ges} = \sin^2\!\bigl(\kappa\sqrt{P_\mathrm{ges}}\bigr) ,
  \qquad
  \eta_n = \frac{P_n}{P_\mathrm{ges}}\;\eta_\mathrm{ges} ,
  \qquad P_\mathrm{ges}=\sum_n P_n .
  \label{eq:multitone}
\end{empheq}
Für $N=1$ ist $P_\mathrm{ges}=P_0$ und Gl.~\eqref{eq:multitone} geht in
Gl.~\eqref{eq:etasingle} über, wie sie muss. Die Energiebilanz stimmt exakt:
$|A_\mathrm{u}|^2+\sum_n|A_n|^2 = \cos^2\kappa_\mathrm{eff}+\sin^2\kappa_\mathrm{eff} = 1$.

\medskip
\noindent\textbf{Die Summe steht also \emph{unter} der Wurzel, nicht außen
vor den Sinusquadraten.} Anschaulich: Die $N$ Gitter wirken wie \emph{ein
einziges} Gitter der Stärke $\sqrt{\sum_n\kappa_n^2}$. Ein zusätzlicher Ton
öffnet keinen neuen Kanal aus dem ungebeugten Strahl heraus -- er benutzt
denselben Kanal mit. Deshalb zählt nur die Gesamtleistung.

\paragraph{Warum die naive Summe trotzdem verbreitet ist.}
Sie ist nicht falsche Physik, sondern die Kleinsignalgrenze, an der falschen
Stelle genommen. Für $\kappa\sqrt{P}\ll1$ ist $\sin^2 x\approx x^2$, und dann
gilt
\begin{equation}
  \sum_n \sin^2\!\bigl(\kappa\sqrt{P_n}\bigr) \approx \kappa^2\sum_n P_n
  = \kappa^2 P_\mathrm{ges} \approx \sin^2\!\bigl(\kappa\sqrt{P_\mathrm{ges}}\bigr) .
\end{equation}
Beide Ausdrücke fallen zusammen. Sie laufen genau dann auseinander, wenn die
Sättigung einsetzt -- und das ist der Bereich, in dem wir arbeiten:

\begin{center}
\begin{tabular}{ccccc}
  \toprule
  $P_\mathrm{ges}$ & $\kappa\sqrt{P_\mathrm{ges}}$ & $\eta_\mathrm{ges}$ korrekt
    & naive Summe (3 Töne) & Fehler\\
  \midrule
  \SI{5}{\milli\watt}   & $0{,}066$ & $0{,}00432$ & $0{,}00433$ & $+0{,}10\,\%$\\
  \SI{50}{\milli\watt}  & $0{,}208$ & $0{,}0427$  & $0{,}0431$  & $+0{,}97\,\%$\\
  \textbf{\SI{0.664}{\watt}} & $0{,}758$ & $\mathbf{0{,}473}$ & $0{,}539$ & $\mathbf{+14{,}0\,\%}$\\
  \bottomrule
\end{tabular}
\end{center}

\noindent Am Arbeitspunkt liegt die naive Summe für die $x$-Achse
\SI{14.0}{\percent} zu hoch, für die $y$-Achse (\SI{0.825}{\watt}, vier Töne)
\SI{20.2}{\percent}, und bei dreizehn Tönen \SI{19.8}{\percent}. Bei
stärkerer Kopplung überschreitet sie sogar \SI{100}{\percent} -- der
sicherste Hinweis, dass sie kein Erhaltungssatz kennt.

\paragraph{Gültigkeitsbereich.}
Gl.~\eqref{eq:multitone} ist \emph{keine} Kleinsignalnäherung, sondern die
exakte Lösung von Gl.~\eqref{eq:coupled} bei beliebiger Kopplungsstärke.
Numerisch gegen eine direkte Integration von Gl.~\eqref{eq:coupled} geprüft:

\begin{center}
\begin{tabular}{cccc}
  \toprule
  $\kappa_\mathrm{eff}$ & $\eta_\mathrm{ges}$ exakt & $\sin^2\kappa_\mathrm{eff}$ & Abweichung\\
  \midrule
  $0{,}211$ & $0{,}044015$ & $0{,}044015$ & $7\cdot10^{-16}$\\
  $0{,}705$ & $0{,}419510$ & $0{,}419510$ & $7\cdot10^{-14}$\\
  $1{,}409$ & $0{,}974086$ & $0{,}974086$ & $5\cdot10^{-13}$\\
  $2{,}114$ & $0{,}733126$ & $0{,}733126$ & $5\cdot10^{-13}$\\
  $2{,}818$ & $0{,}100972$ & $0{,}100972$ & $1\cdot10^{-13}$\\
  \bottomrule
\end{tabular}
\end{center}

\noindent Die letzten drei Zeilen liegen jenseits des Maximums, wo die
Effizienz durch Rückkonversion wieder fällt -- auch dort stimmt die Formel,
und die Aufteilung $\eta_n/\eta_\mathrm{ges} = P_n/P_\mathrm{ges}$ bleibt auf
sechs Stellen konstant.

Vorausgesetzt wird stattdessen dreierlei, und das sind die eigentlichen
Grenzen.

\medskip
\noindent\textbf{(i) Gemeinsame Phasenanpassung.} Alle $N$ Gitter müssen
dieselbe Bragg-Bedingung erfüllen. Sie tun das nicht exakt: Zu jedem Ton
gehört ein leicht anderer Bragg-Winkel, und bei festem Einfall bleibt ein
Impulsfehler $\Delta k$. Die verstimmte Zwei-Wellen-Lösung lautet
\begin{equation}
  \eta(\xi) = \frac{\nu^2}{\nu^2+\xi^2}\,\sin^2\!\sqrt{\nu^2+\xi^2} ,
  \qquad \nu = \kappa\sqrt{P},\quad \xi = \frac{\Delta k\,L}{2} .
  \label{eq:detuned}
\end{equation}
Die Verstimmung $\xi$ muss nicht aus der Geometrie gerechnet werden -- sie
steckt bereits in der spezifizierten Bandbreite. Definitionsgemäß ist die
\SI{36}{\mega\hertz} genau die Spanne, über der die Effizienz auf die Hälfte
fällt; am Bandrand $\delta f = \pm\SI{18}{\mega\hertz}$ gilt also
$\eta=\tfrac12\eta(0)$. Bei $\nu=\kappa\sqrt{0{,}664}=0{,}758$ löst
Gl.~\eqref{eq:detuned} das zu $\xi_{1/2}=1{,}3625$, und weil $\xi$ linear in
der Frequenzablage ist, folgt die Eichung
\begin{equation}
  \xi(\delta f) = 1{,}3625\cdot\frac{\delta f}{\SI{18}{\mega\hertz}} .
\end{equation}
Damit lässt sich jede Tonspanne beziffern -- $\delta f$ ist die Ablage des
äußersten Tons, also die halbe Spanne:

\begin{center}
\begin{tabular}{cccc}
  \toprule
  Tonspanne & $\delta f$ (äußerster Ton) & $\xi$ & Effizienzverlust dort\\
  \midrule
  \textbf{\SI{0.37}{\mega\hertz}} & \SI{0.185}{\mega\hertz} & $0{,}014$ & $\mathbf{0{,}007\,\%}$\\
  \SI{6}{\mega\hertz}  & \SI{3}{\mega\hertz}  & $0{,}227$ & $1{,}8\,\%$\\
  \SI{18}{\mega\hertz} & \SI{9}{\mega\hertz}  & $0{,}681$ & $15{,}1\,\%$\\
  \SI{36}{\mega\hertz} & \SI{18}{\mega\hertz} & $1{,}363$ & $50\,\%$ (Bandrand, Eichpunkt)\\
  \bottomrule
\end{tabular}
\end{center}

\noindent Die Tonspanne des Arbeitspunkts ist \emph{ein Hundertstel Prozent}
der Bandbreite wert -- die Phasenanpassung ist hier schlicht kein Thema. Das
ist keine Selbstverständlichkeit: Ein Profil mit \SI{18}{\mega\hertz} Spanne,
wie es größere Spotfelder brauchen, verlöre am Rand bereits
\SI{15}{\percent}, und dann wäre Gl.~\eqref{eq:multitone} durch die
verstimmte Version Gl.~\eqref{eq:detuned} je Ton zu ersetzen, mit
ordnungsabhängigem $\eta_n$.

\medskip
\noindent\textbf{(ii) Bragg-Regime, zwei Wellen je Gitter.} Höhere
Beugungsordnungen und optische Mischprodukte sind stark fehlangepasst und
vernachlässigt.

\medskip
\noindent\textbf{(iii) Linearität des Wandlers.} $\kappa_n\propto\sqrt{P_n}$
setzt voraus, dass die HF-Leistung linear in die Indexmodulation übergeht.
Akustische Nichtlinearität ist nicht enthalten -- dazu gleich mehr, denn an
dieser Stelle kommt der Crestfaktor ein zweites Mal ins Spiel.

\paragraph{Warum die Tonphasen hier nirgends auftauchen.}
Das ist kein Versehen, sondern der Kern der Abgrenzung aus
Abschnitt~\ref{sec:abgrenzung}. Die akustooptische Wechselwirkung sieht die
\emph{Fourierkomponenten} des Schallfeldes. Jeder Ton ist genau eine solche
Komponente, mit zeitlich konstanter Amplitude $a_n$ und konstanter Frequenz;
er bildet ein stehendes Gitter, dessen Beugung von $P_n = a_n^2$ abhängt und
von sonst nichts. Die Phase $\varphi_n$ legt nur fest, \emph{wo} die
Gitterberge dieses einen Tons liegen, und das überträgt sich als derselbe
Phasenversatz auf die gebeugte Welle -- also auf $\varphi_{nm}$ im
Interferenzbild, wo er hingehört, und nicht auf $\eta_n$.

Der Crestfaktor dagegen ist eine Eigenschaft der \emph{Summe} $s(t)$, nicht
einer einzelnen Komponente. Diese Summe existiert nur dort, wo alle Töne
durch dasselbe breitbandige Bauteil müssen: im Verstärker. Damit ist
endgültig klar, warum der Crestfaktor eine Budgetgrenze der Elektronik ist
und in der ganzen Beugungsrechnung nicht vorkommt.

\medskip
\noindent\emph{Mit einer Ausnahme}, und das ist Punkt (iii): Die momentane
Dehnung an einem Ort im Kristall ist $s(t)$, mit Crestfaktor $\crest$. Wird
die Spitzendehnung groß genug, erzeugt der Kristall selbst Harmonische, und
die Zerlegung in $N$ unabhängige Gitter gilt nicht mehr. Das ist ein
zweiter, vom Verstärker unabhängiger Kanal, über den ein hoher Crestfaktor
schadet -- und er zeigt in dieselbe Richtung. Quantifizieren lässt er sich
hier nicht; dafür bräuchte es eine Messung der Oberwellen bei fester
mittlerer Leistung und variiertem $\crest$.

\paragraph{Die Merkregel.}
Aus Gl.~\eqref{eq:multitone} folgt unmittelbar eine Aussage, die man sich
merken sollte:

\begin{center}
\emph{Die Gesamtbeugungseffizienz hängt nur von der gesamten HF-Leistung ab,
nicht davon, auf wie viele Töne sie verteilt wird.}
\end{center}

\noindent Bei \SI{0.664}{\watt} ist $\eta_\mathrm{ges}=0{,}473$, ob die
Leistung auf einen, drei oder dreizehn Töne verteilt wird; nur die Aufteilung
auf die Ordnungen ändert sich (je Ton $0{,}473$, $0{,}158$ bzw.\ $0{,}036$).
Das Aufspalten in ein Multitonprofil kostet also \emph{keine} Effizienz -- die
geringere Effizienz gegenüber dem Testsheet kommt allein aus der kleineren
Gesamtleistung (\SI{0.66}{\watt} statt \SI{1.3}{\watt}).

\paragraph{Die Zahlen.}
Mit $P_\mathrm{mittel,x}=\SI{0.664}{\watt}$ und
$P_\mathrm{mittel,y}=\SI{0.825}{\watt}$ aus Schritt~1:
\begin{align}
  \eta_x &= \sin^2\!\bigl(0{,}930\cdot\sqrt{0{,}664}\bigr) = 0{,}473 ,\\
  \eta_y &= \sin^2\!\bigl(0{,}930\cdot\sqrt{0{,}825}\bigr) = 0{,}559 .
\end{align}
Die Leistungsanteile aus $r_x=0{,}97$ und $r_y=1{,}16$ verteilen das nach
Gl.~\eqref{eq:multitone} auf die Ordnungen,
\begin{align}
  \eta_{x,n} &= [\,0{,}156;\;0{,}161;\;0{,}156\,], \\
  \eta_{y,m} &= [\,0{,}150;\;0{,}130;\;0{,}130;\;0{,}150\,],
\end{align}
und weil jeder Spot \emph{zweimal} gebeugt wird -- einmal im $x$-, einmal im
$y$-AOD --, ist die Effizienz des Spots $(n,m)$ das Produkt
$\eta_{x,n}\,\eta_{y,m}$. Das erklärt zugleich, wo das ungebeugte Licht
bleibt: Der in $x$ nicht gebeugte Anteil $1-\eta_x$ durchläuft den
$y$-AOD weiter und landet in den Ordnungen „0-$m$``, der in $y$ nicht
gebeugte in „$n$-0``, und $(1-\eta_x)(1-\eta_y)$ geht ganz geradeaus
als „0-0``. Diese Strahlen laufen geometrisch getrennt und werden
ausgeblendet -- deshalb spricht das Testsheet von der Ordnung „1-1``, und
deshalb ist nur sie das Nutzsignal. Summiert über alle zwölf Spots faktorisiert die
Doppelsumme:
\begin{equation}
  \eta_\mathrm{Profil} \;=\; \sum_{n,m}\eta_{x,n}\,\eta_{y,m}
  \;=\; \Bigl(\sum_n \eta_{x,n}\Bigr)\Bigl(\sum_m \eta_{y,m}\Bigr)
  \;=\; \eta_x\,\eta_y \;=\; \SI{26.4}{\percent},
\end{equation}
also zwischen \SI{2.0}{\percent} und \SI{2.4}{\percent} je Spot, je nach
Randgewichtung.

\paragraph{Intermodulation.}
Mit $\mathrm{OIP3}=+\SI{49}{dBm}$ und der stärksten Tonleistung
$\SI{0.226}{\watt} = +\SI{23.5}{dBm}$ liegen die Produkte dritter Ordnung
\begin{equation}
  \mathrm{IMD3} = 2\,\bigl(\mathrm{OIP3} - P_\mathrm{Ton}\bigr)
  = 2\,(49 - 23{,}5) = \SI{51}{dBc}
\end{equation}
unter den Haupttönen. Geisterspots \SI{51}{\deci\bel} unter den Hauptspots
sind unkritisch. Bei $P_\mathrm{mittel}=\SI{1.3}{\watt}$ wären es noch
\SI{45}{dBc} -- dort gilt die Extrapolation über OIP3 allerdings ohnehin nicht
mehr, weil der Verstärker in der Begrenzung sitzt.

\subsection{Schritt 4: Was die Phasenwahl in Watt bringt}

Jetzt zahlt sich alles Vorangegangene aus. Nach Gl.~\eqref{eq:pavglimit}
übersetzt sich jeder Crestfaktor unmittelbar in erlaubte Leistung und damit in
Licht:

\begin{table}[htb]
  \centering
  \caption{Alle Sätze bei $P_\mathrm{peak}\le P_{1\mathrm{dB}}$. Die
  $x$-Achse liegt bei jedem zentrierten Satz fest bei $\crest_x=2{,}449$.}
  \label{tab:hardware}
  \begin{tabular}{lcccc}
    \toprule
    Phasensatz & $\crest_y$ & $P_\mathrm{mittel,y}$ & $\eta_\mathrm{Profil}$ & Versatz / $\Uw$\\
    \midrule
    alle Phasen $0$      & $2{,}826$ & \SI{0.50}{\watt} & \SI{17.6}{\percent} & \SI{0}{\nano\meter} / \SI{0.18}{\percent}\\
    frühere Wahl         & $2{,}485$ & \SI{0.65}{\watt} & \SI{21.8}{\percent} & \SI{0}{\nano\meter} / \SI{0.21}{\percent}\\
    \textbf{gewählter Punkt} & $\mathbf{2{,}196}$ & $\mathbf{\SI{0.83}{\watt}}$ & $\mathbf{\SI{26.4}{\percent}}$ & \SI{0}{\nano\meter} / \SI{0.30}{\percent}\\
    Schroeder            & $2{,}000$ & \SI{1.00}{\watt} & \SI{36.8}{\percent} & \SI{18.9}{\nano\meter} / \SI{17.9}{\percent}\\
    \bottomrule
  \end{tabular}
\end{table}

Der gewählte Satz nach Gl.~\eqref{eq:workingpoint} schöpft die einzige
verbliebene Freiheit aus: \SI{0.83}{\watt} statt \SI{0.50}{\watt} gegenüber
$\varphi=0$ und \SI{0.65}{\watt} gegenüber der früheren Wahl, bei unverändert
exakter Zentrierung und $\Uw=\SI{0.30}{\percent}$. In Licht sind das
\SI{26.4}{\percent} statt \SI{17.6}{\percent} bzw.\ \SI{21.8}{\percent} im
Profil.

Schroeder brächte mit \SI{36.8}{\percent} nochmals mehr, verletzt aber
Gl.~\eqref{eq:centering}: Der Fleck sitzt \SI{18.9}{\nano\meter} neben dem
Atom und $\Uw$ springt um zwei Größenordnungen. Das ist der bereits bekannte
Preis der Zentrierung, hier in echten Watt.

\section*{Offene Punkte}

\begin{itemize}\itemsep0.2em
  \item Die Datenblattwerte des Verstärkers sind übernommen, die
        Kompressionskurve der konkreten Kette ist \emph{nicht gemessen}.
        Gl.~\eqref{eq:powerbudget} unterstellt eine harte Spitzenspannungs\-grenze
        mit linearem Verhalten bis dorthin; die reale Kompressionskurve,
        die Sättigung des AOD selbst (etwa $\eta_\mathrm{diff}/3$), die
        akustische Dämpfung sowie Filter, Kabel und Wandler sind nicht
        enthalten. Der berechnete Wert gilt für das ideale AWG-Signal.
  \item Bei $N_x=3$ ist die Spitzenlast der $x$-Achse durch
        Gl.~\eqref{eq:centering} auf $\sqrt{6}$ festgelegt. Eine Entlastung
        dieser Achse ist nur über mehr Töne oder über das Aufgeben der
        exakten Zentrierung erreichbar.
  \item Der Atomort ist fest der Profilmittelpunkt; Strahllagefluktuationen
        sind nicht berücksichtigt.
\end{itemize}

\end{document}
